% ===================================================================
% Trajectory Completion Computing for Autonomous Vehicles:
% A Categorical Architecture Replacing Forward Simulation
%
% Author: Kundai Farai Sachikonye
% Technical University of Munich
% ===================================================================

\documentclass[twocolumn,10pt]{article}

% --- Packages ---
\usepackage[utf8]{inputenc}
\usepackage[T1]{fontenc}
\usepackage{amsmath,amssymb,amsthm,amsfonts}
\usepackage{mathtools}
\usepackage{geometry}
\usepackage{graphicx}
\usepackage{float}
\usepackage{algorithm}
\usepackage{algorithmic}
\usepackage{booktabs}
\usepackage{array}
\usepackage{hyperref}
\usepackage{cleveref}
\usepackage{enumitem}
\usepackage{caption}
\usepackage{subcaption}
\usepackage{xcolor}
\usepackage{microtype}
\usepackage{fancyhdr}
\usepackage{titlesec}

\geometry{margin=0.75in,top=1in,bottom=1in}

% --- Theorem environments ---
\newtheorem{theorem}{Theorem}[section]
\newtheorem{lemma}[theorem]{Lemma}
\newtheorem{proposition}[theorem]{Proposition}
\newtheorem{corollary}[theorem]{Corollary}
\theoremstyle{definition}
\newtheorem{definition}[theorem]{Definition}
\newtheorem{example}[theorem]{Example}
\newtheorem{remark}[theorem]{Remark}
\newtheorem{axiom}[theorem]{Axiom}

% --- Custom commands ---
\newcommand{\Sentropy}{\mathbf{S}}
\newcommand{\Sk}{S_k}
\newcommand{\St}{S_t}
\newcommand{\Se}{S_e}
\newcommand{\dcat}{d_{\mathrm{cat}}}
\newcommand{\Ocat}{\hat{O}_{\mathrm{cat}}}
\newcommand{\Ophys}{\hat{O}_{\mathrm{phys}}}
\newcommand{\Cn}{C(n)}
\newcommand{\taup}{\tau_p}
\newcommand{\Scurrent}{S_{\mathrm{current}}}
\newcommand{\Snext}{S_{\mathrm{next}}}
\newcommand{\Sfinal}{S_{\mathrm{final}}}
\newcommand{\Spen}{S_{\mathrm{pen}}}
\newcommand{\BigO}{\mathcal{O}}

% --- Title ---
\title{\textbf{Trajectory Completion Computing for Autonomous Vehicles:\\A Categorical Architecture Replacing Forward Simulation}}

\author{Kundai Farai Sachikonye\\
\textit{Technical University of Munich}\\
\texttt{kundai.sachikonye@wzw.tum.de}}

\date{}

\begin{document}

\maketitle

% ===================================================================
\begin{abstract}
Every autonomous vehicle architecture deployed or proposed to date implements the same fundamental pattern: perceive the environment, predict future states, plan a trajectory through predicted states, and execute control commands to follow that trajectory. We prove that this pipeline is \emph{architecturally incapable} of safe autonomous driving. The failure is not engineering immaturity but mathematical impossibility: Lyapunov divergence renders prediction useless in $\BigO(1/\lambda)$ seconds, and information-theoretic analysis shows prediction requires $10^3$ to $10^7$ times more bits than driving actually consumes. We present an alternative architecture grounded in the \emph{trajectory completion} framework, where observation, computation, and processing are identified as a single categorical operation: $O(x) \equiv C(x) \equiv P(x)$. The vehicle is modeled as a network of coupled counting loops---oscillators whose tick accumulation \emph{simultaneously} constitutes observation, computation, and actuation. Navigation proceeds by \emph{backward} trajectory completion in $\BigO(\log_3 N)$ time, and driving decisions reduce to \emph{sufficiency recognition} via triple convergence of thermodynamic, categorical, and partition perspectives. We prove partition boundedness ($\dcat \leq \sqrt{3}$ for all time), G\"{o}delian safety ($\varepsilon > 0$ ensures recognition), and thermodynamic security (unauthorized transitions produce detectable heating). The architecture requires no prediction, no simulation, no separate planning module, and no distinction between computer and car. We formally specify five subsystems---Categorical State Manager, Penultimate Navigation Engine, Sufficiency Recognition Module, Completion Morphism Executor, and Triple Equivalence Monitor---and prove safety guarantees for each. Inter-vehicle coordination emerges as a thermodynamic phase transition in a categorical gas, requiring no explicit vehicle-to-vehicle communication protocol.
\end{abstract}

\noindent\textbf{Keywords:} autonomous driving, trajectory completion, categorical computing, coupled oscillators, S-entropy, backward navigation, sufficiency recognition, thermodynamic security

% ===================================================================
\section{Introduction}
\label{sec:introduction}

The autonomous vehicle industry has invested over \$100 billion in a computational architecture whose mathematical foundations guarantee its failure. Every deployed system---from Waymo's modular stack~\cite{urmson2008,levinson2011} to Tesla's end-to-end neural approach~\cite{bojarski2016,chen2024} to Cruise's hybrid pipeline~\cite{bansal2019}---implements a common pattern:
\begin{equation}
\text{Perceive} \to \text{Predict} \to \text{Plan} \to \text{Control}
\label{eq:conventional-pipeline}
\end{equation}
Each stage processes the output of its predecessor and feeds its successor. The architecture appears natural: one must see the world, anticipate what will happen, decide what to do, and then do it.

We argue this architecture is not merely difficult to perfect but \emph{mathematically impossible} to make safe. The argument proceeds in three steps.

\textbf{First}, prediction of chaotic systems diverges exponentially~\cite{poincare1890}. Traffic is a many-body system with Lyapunov exponents $\lambda > 0$, so prediction error $\varepsilon(t) \sim \varepsilon_0 e^{\lambda t}$ renders any forecast useless in $\BigO(1/\lambda)$ seconds. For highway driving, $\lambda \approx 0.5\text{--}2.0\;\text{s}^{-1}$, giving useful prediction horizons of $0.5$--$2$ seconds---far shorter than the $5$--$15$ second horizons required for safe driving~\cite{schwarting2018}.

\textbf{Second}, the information content required for prediction vastly exceeds the information consumed by driving. Let $H(D)$ denote the entropy of the full driving environment and $I(D; A_Q)$ the mutual information between the environment and the action actually taken at categorical resolution $Q$. We show that the ratio $I(D; A_Q)/H(D) \sim 10^{-3}$ to $10^{-7}$. Prediction-based architectures must process $H(D)$; trajectory completion architectures process only $I(D; A_Q)$.

\textbf{Third}, the pipeline architecture~\eqref{eq:conventional-pipeline} is \emph{serially fragile}: failure at any stage propagates to all downstream stages. There is no feedback mechanism that can compensate faster than the pipeline latency. A perception failure (missed object) produces a prediction failure (absent agent), which produces a planning failure (trajectory through occupied space), which produces a control failure (collision). Each of the documented fatal AV accidents traces to this serial fragility.

We propose a fundamentally different architecture based on \emph{trajectory completion computing}~\cite{sachikonye2026trajectory,sachikonye2026backward}. The foundational identity is:
\begin{equation}
\boxed{O(x) \equiv C(x) \equiv P(x)}
\label{eq:fundamental-identity}
\end{equation}
where $O$ denotes observation, $C$ denotes computation, and $P$ denotes physical processing. All three are the same operation: \emph{categorical address resolution in bounded phase space}. This identity eliminates the pipeline entirely. There is no perception stage (observation \emph{is} computation), no prediction stage (backward completion replaces forward simulation), no planning stage (the penultimate state is unique), and no control stage (the completion morphism \emph{is} the physical transition).

The vehicle is not a computer mounted on a car. It is a \emph{single coupled oscillator network} in which steering columns, engine pistons, brake rotors, wheel bearings, sensor elements, and processor clock crystals are all oscillators in a shared partition space. The ``driving decision'' is not computed and then executed---it \emph{emerges} from the coupled dynamics of counting loops.

This paper makes the following contributions:
\begin{enumerate}[leftmargin=*]
\item We prove that forward simulation architectures have unbounded failure probability beyond a finite horizon (\Cref{sec:forward-fails}).
\item We define the \emph{counting loop} as the primitive computational element and prove it simultaneously performs observation, computation, and processing (\Cref{sec:counting-loop}).
\item We specify a complete five-subsystem architecture with formal interface contracts (\Cref{sec:architecture}--\Cref{sec:tem}).
\item We prove partition-bounded safety, G\"{o}delian recognition guarantees, and thermodynamic security (\Cref{sec:safety}).
\item We derive inter-vehicle coordination as a thermodynamic phase transition requiring no explicit communication (\Cref{sec:coordination}).
\end{enumerate}

The paper is organized as follows. \Cref{sec:forward-fails} provides a rigorous analysis of why forward simulation fails. \Cref{sec:counting-loop} develops the counting loop primitive. \Cref{sec:architecture} presents the architectural overview. \Cref{sec:csm}--\Cref{sec:tem} detail each subsystem. \Cref{sec:safety} proves safety guarantees. \Cref{sec:coordination} treats inter-vehicle coordination. \Cref{sec:comparison} compares with conventional architectures. \Cref{sec:implementation} discusses implementation. \Cref{sec:conclusion} concludes.


% ===================================================================
\section{Why Forward Simulation Fails}
\label{sec:forward-fails}

We now make precise the claim that the conventional autonomous driving pipeline is architecturally incapable of achieving safety. The argument is threefold: computational complexity, dynamical instability, and information-theoretic waste.

\subsection{Computational Complexity of Forward Simulation}

\begin{definition}[Forward Simulation]
\label{def:forward-sim}
A \emph{forward simulation} of $n$ interacting agents over time horizon $T$ with time step $\Delta t$ computes, for each agent $i$ and each time step $k$:
\begin{equation}
\mathbf{x}_i(k+1) = f\bigl(\mathbf{x}_1(k), \ldots, \mathbf{x}_n(k), \mathbf{u}_i(k)\bigr)
\end{equation}
where $\mathbf{x}_i \in \mathbb{R}^d$ is the state of agent $i$ and $\mathbf{u}_i$ is the control input.
\end{definition}

\begin{theorem}[Forward Simulation Complexity]
\label{thm:forward-complexity}
The computational complexity of forward simulation for $n$ interacting agents over time horizon $T$ with time step $\Delta t$ is
\begin{equation}
\mathcal{C}_{\mathrm{fwd}} = \BigO\!\left(\frac{n^2 T}{\Delta t}\right).
\end{equation}
\end{theorem}

\begin{proof}
At each of the $T/\Delta t$ time steps, each agent's state update depends on all $n$ agents (pairwise interaction), requiring $\BigO(n^2)$ evaluations per step. The product gives the stated complexity. Even with spatial partitioning reducing pairwise checks to $\BigO(n \log n)$ per step, the $T/\Delta t$ factor remains, and the fundamental scaling is superlinear in both $n$ and $T/\Delta t$.
\end{proof}

In a typical urban environment, $n \approx 50$--$200$ agents, $T \approx 10$ s, and $\Delta t \approx 0.01$ s (for accurate physics), yielding $\mathcal{C}_{\mathrm{fwd}} \approx 10^7$--$10^9$ operations per planning cycle. This is feasible on modern hardware, but the \emph{accuracy} of this computation is the fatal issue, as we now show.

\subsection{Lyapunov Divergence}

\begin{theorem}[Prediction Divergence]
\label{thm:lyapunov}
Let $\varepsilon_0$ be the initial state estimation error and $\lambda_{\max} > 0$ the maximal Lyapunov exponent of the traffic system. Then the prediction error at time $t$ satisfies
\begin{equation}
\varepsilon(t) \geq \varepsilon_0 \, e^{\lambda_{\max} t}.
\label{eq:lyapunov-divergence}
\end{equation}
For any safety threshold $\varepsilon_{\mathrm{safe}}$, the prediction becomes unsafe at
\begin{equation}
T_{\mathrm{safe}} = \frac{1}{\lambda_{\max}} \ln\!\left(\frac{\varepsilon_{\mathrm{safe}}}{\varepsilon_0}\right).
\end{equation}
\end{theorem}

\begin{proof}
By definition, the maximal Lyapunov exponent satisfies
\[
\lambda_{\max} = \lim_{t \to \infty} \frac{1}{t} \ln \frac{\|\delta \mathbf{x}(t)\|}{\|\delta \mathbf{x}(0)\|}
\]
for generic perturbations $\delta \mathbf{x}(0)$~\cite{arnold1978}. For finite $t$, we have $\|\delta \mathbf{x}(t)\| \geq \|\delta \mathbf{x}(0)\| e^{\lambda_{\max} t}$ along the most unstable direction. Setting $\|\delta \mathbf{x}(0)\| = \varepsilon_0$ and solving $\varepsilon_0 e^{\lambda_{\max} T_{\mathrm{safe}}} = \varepsilon_{\mathrm{safe}}$ yields the stated bound.
\end{proof}

\begin{corollary}
\label{cor:prediction-horizon}
With $\varepsilon_0 = 0.1$ m (sensor accuracy), $\varepsilon_{\mathrm{safe}} = 1.0$ m (lane width margin), and $\lambda_{\max} = 1.0\;\mathrm{s}^{-1}$ (moderate traffic), we obtain $T_{\mathrm{safe}} = \ln(10) \approx 2.3$ seconds. Any prediction beyond $2.3$ s is \emph{mathematically guaranteed} to violate safety bounds.
\end{corollary}

This result is not an artifact of poor modeling. It is a consequence of Poincar\'{e}'s foundational work on the three-body problem~\cite{poincare1890}: systems with more than two interacting bodies exhibit chaotic sensitivity to initial conditions. Traffic, with dozens to hundreds of interacting agents, each possessing autonomous decision-making capability, is deeply in this regime.

\subsection{The Prediction Paradox}

\begin{remark}[The Prediction Paradox]
\label{rem:prediction-paradox}
Humans cannot predict traffic either. A human driver cannot state where each surrounding vehicle will be in 10 seconds. Yet humans drive safely---billions of person-hours per year with remarkably low accident rates. This establishes an existence proof: safe driving does \emph{not require prediction}. The conventional AV architecture demands a capability (accurate long-horizon prediction) that is (a)~mathematically impossible and (b)~empirically unnecessary.
\end{remark}

What humans do instead is \emph{recognize sufficiency}: they assess whether the current state affords enough margin for the next action. This assessment is local in time (no prediction horizon needed), categorical in nature (not metric---``enough space'' rather than ``3.7 meters''), and multi-perspective (visual, auditory, proprioceptive agreement). Our architecture formalizes precisely this capability.

The prediction paradox has a deeper implication. Every AV accident report in the public record identifies the same root cause pattern: the system's prediction of another agent's behavior was wrong. The pedestrian stepped off the curb unexpectedly. The cyclist swerved. The truck braked suddenly. In each case, the system had predicted a different trajectory and planned around that prediction. When the prediction failed, the plan was invalid, and the control system executed an invalid plan. The architecture made the accident \emph{inevitable} once the prediction failed, and \Cref{thm:lyapunov} guarantees that predictions fail.

\subsection{Information-Theoretic Analysis}

\begin{theorem}[Information-Theoretic Waste of Prediction]
\label{thm:info-waste}
Let $H(D)$ denote the Shannon entropy of the driving environment at sensor resolution, and let $I(D; A_Q)$ denote the mutual information between the environment and the driving action at categorical resolution $Q$. Then
\begin{equation}
\frac{I(D; A_Q)}{H(D)} \in [10^{-7}, 10^{-3}].
\end{equation}
\end{theorem}

\begin{proof}
The driving environment entropy comprises: LiDAR point clouds ($\sim 10^6$ points $\times$ 16 bits per coordinate $\times$ 3 coordinates $= 4.8 \times 10^7$ bits/frame), cameras ($6 \times 2 \times 10^6$ pixels $\times$ 24 bits $\approx 2.88 \times 10^8$ bits/frame), radar ($\sim 10^4$ detections $\times$ 64 bits $\approx 6.4 \times 10^5$ bits/frame), IMU ($6$ axes $\times$ 32 bits $= 192$ bits/frame), GPS ($3 \times 64 = 192$ bits/frame), and map data ($\sim 10^6$ bits cached). Summing: $H(D) \approx 3.4 \times 10^8$ bits/frame at the lower end, up to $\sim 10^9$ bits/frame with high-resolution sensors~\cite{shannon1948,cover2006}.

The driving action at categorical resolution $Q$ selects among a finite set of transitions. The action space decomposes as: longitudinal $\in \{$accelerate strongly, accelerate mildly, maintain, decelerate mildly, decelerate strongly$\}$ and lateral $\in \{$turn left sharply, turn left mildly, straight, turn right mildly, turn right sharply$\}$. This gives $|A_Q| \leq 5 \times 5 = 25$ combined actions, so $H(A_Q) \leq \log_2 25 \approx 4.6$ bits. The mutual information satisfies $I(D; A_Q) \leq H(A_Q) \leq 4.6$ bits. In practice, the conditional entropy $H(A_Q | D)$ is small (the action is nearly determined by the state), so $I(D; A_Q) \approx H(A_Q) - H(A_Q | D) \approx 2$--$5$ bits per decision cycle.

The ratio is therefore $2/10^9 \leq I(D; A_Q)/H(D) \leq 5/10^5$, which lies in $[10^{-7}, 10^{-3}]$ depending on sensor suite and resolution.
\end{proof}

This theorem reveals the colossal waste of the prediction paradigm. A prediction-based system must process the full environment entropy $H(D) \sim 10^8$--$10^9$ bits in order to extract the $2$--$5$ bits of mutual information actually needed for driving. The trajectory completion architecture processes only the $I(D; A_Q)$ bits directly, achieving a compression ratio of $10^3$ to $10^7$.

\begin{theorem}[Unbounded Failure Probability]
\label{thm:unbounded-failure}
Any prediction-based architecture has failure probability $P_{\mathrm{fail}}(t) \to 1$ as $t \to \infty$, regardless of the quality of the prediction model.
\end{theorem}

\begin{proof}
By \Cref{thm:lyapunov}, for any finite prediction horizon $T_h$, there exists a time $T_{\mathrm{safe}} < T_h$ beyond which the prediction error exceeds safety bounds. The probability that the system encounters a scenario requiring decisions beyond $T_{\mathrm{safe}}$ is nonzero for any nontrivial driving environment---such scenarios include merging, intersection negotiation, pedestrian crossing, and emergency maneuvers. Let $p_{\mathrm{encounter}}$ denote the per-unit-time probability of encountering such a scenario, with $p_{\mathrm{encounter}} > 0$.

Over time $t$, the expected number of critical encounters is $N(t) = p_{\mathrm{encounter}} \cdot t \to \infty$. Each encounter has a nonzero probability of failure $p_k \geq p_{\min} > 0$ because the prediction is unreliable beyond $T_{\mathrm{safe}}$. The cumulative survival probability is
\[
P_{\mathrm{survive}}(t) = \prod_{k=1}^{N(t)} (1 - p_k) \leq (1 - p_{\min})^{N(t)} \to 0
\]
as $N(t) \to \infty$. Therefore $P_{\mathrm{fail}}(t) = 1 - P_{\mathrm{survive}}(t) \to 1$.
\end{proof}

This theorem explains why every autonomous vehicle company accumulates accidents with deployment time. The failures are not bugs to be fixed---they are theorems to be acknowledged. The architecture itself is the problem.

\begin{lemma}[Pipeline Latency Amplification]
\label{lem:pipeline-latency}
In the sequential pipeline~\eqref{eq:conventional-pipeline}, the total latency is
\begin{equation}
\tau_{\mathrm{total}} = \tau_{\mathrm{perceive}} + \tau_{\mathrm{predict}} + \tau_{\mathrm{plan}} + \tau_{\mathrm{control}} \geq 4 \tau_{\min}
\end{equation}
where $\tau_{\min}$ is the minimum single-stage latency. During $\tau_{\mathrm{total}}$, the environment has evolved, rendering the perception data stale.
\end{lemma}

\begin{proof}
Each stage requires at least $\tau_{\min}$ to execute (bounded below by data transfer, computation, and actuation latencies). The stages are serial: prediction cannot begin until perception completes, planning cannot begin until prediction completes, and control cannot begin until planning completes. The sum is therefore at least $4\tau_{\min}$. With typical values $\tau_{\min} \approx 10$--$50$ ms per stage, the total latency is $40$--$200$ ms. At highway speeds ($30$ m/s), the vehicle travels $1.2$--$6$ m during this latency---well within the prediction error bound of \Cref{cor:prediction-horizon}. The perception data used for the current control action was acquired $40$--$200$ ms ago, describing a world that no longer exists.
\end{proof}


% ===================================================================
\section{The Counting Loop Primitive}
\label{sec:counting-loop}

Having established that forward simulation is fundamentally inadequate, we now develop the computational primitive that replaces it.

\subsection{Bounded Phase Space}

\begin{axiom}[Bounded Phase Space]
\label{ax:bounded}
The phase space $\Gamma$ of any physical system is bounded: $\Gamma \subset B_R \subset \mathbb{R}^{2d}$ for some finite $R > 0$.
\end{axiom}

This axiom is physically motivated: energy is finite (the vehicle has a finite fuel/battery capacity), momentum is bounded by $E/c$ (relativistically, and practically by engine power), and spatial extent is finite for any real system (the road network is finite)~\cite{arnold1978}. For an automobile on Earth's road network, the phase space is bounded by maximum velocity ($\sim 70$ m/s), maximum acceleration ($\sim 10$ m/s$^2$), road network extent ($\sim 10^7$ m), and energy capacity ($\sim 10^9$ J for a full tank).

Boundedness is the critical structural property that enables trajectory completion. In an unbounded phase space, states can diverge without limit (\Cref{thm:lyapunov}). In a bounded phase space, Poincar\'{e}'s recurrence theorem~\cite{poincare1890} guarantees that the system returns arbitrarily close to any previous state. This transforms the problem from extrapolation (forward prediction in unbounded space) to recognition (identifying which bounded state the system currently occupies).

\begin{definition}[S-entropy Triple]
\label{def:s-entropy}
The \emph{S-entropy} of a system state is the triple
\begin{equation}
\Sentropy = (\Sk, \St, \Se) \in [0,1]^3
\end{equation}
where $\Sk$ is the \emph{categorical} (structural/spatial) entropy measuring the complexity of the system's spatial configuration, $\St$ is the \emph{temporal} (dynamical) entropy measuring the rate of state change, and $\Se$ is the \emph{energetic} (thermodynamic) entropy measuring the degree of energy dispersal, each normalized to $[0,1]$~\cite{sachikonye2026trajectory}.
\end{definition}

For a vehicle, the S-entropy triple has intuitive meaning: $\Sk$ encodes where the vehicle is (lane position, road identity, surrounding geometry), $\St$ encodes how fast things are changing (speed, acceleration, turn rate), and $\Se$ encodes the thermodynamic state (engine temperature, tire temperature, fuel level).

\begin{definition}[Partition Coordinates]
\label{def:partition-coords}
A state in bounded phase space is addressed by \emph{partition coordinates} $(n, \ell, m, s)$ where $n \in \mathbb{N}^+$ is the principal depth (resolution level), $\ell$ is the angular momentum quantum number ($0 \leq \ell < n$), $m$ is the magnetic quantum number ($-\ell \leq m \leq \ell$), and $s \in \{-1/2, +1/2\}$ is the spin. The number of categorical cells at depth $n$ is
\begin{equation}
\Cn = 2n^2.
\label{eq:degeneracy}
\end{equation}
\end{definition}

The partition coordinates provide a \emph{hierarchical address system} for states in phase space. At depth $n=1$, there are $C(1) = 2$ cells (coarsest binary resolution: e.g., ``moving'' vs ``stopped''). At depth $n=10$, there are $C(10) = 200$ cells (fine resolution: specific lane position, speed range, energy state). The structure is a ternary tree: each cell at depth $n$ is the parent of at most three children at depth $n+1$~\cite{sachikonye2026poincare}.

\subsection{The Counting Loop}

\begin{definition}[Counting Loop]
\label{def:counting-loop}
A \emph{counting loop} is a triple $\mathcal{L} = (\omega, \mathcal{P}, \tau)$ where:
\begin{itemize}[leftmargin=*]
\item $\omega \in \mathbb{R}^+$ is an \emph{oscillator} with natural frequency $f = \omega / 2\pi$,
\item $\mathcal{P}$ is a \emph{partition cell} (a bounded region of phase space addressed by coordinates $(n, \ell, m, s)$),
\item $\tau: \mathbb{N} \to \mathcal{P}$ is a \emph{tick map} that associates each oscillator period with a categorical state count in $\mathcal{P}$.
\end{itemize}
\end{definition}

\begin{definition}[Tick]
\label{def:tick}
A single \emph{tick} of counting loop $\mathcal{L} = (\omega, \mathcal{P}, \tau)$ is the atomic operation:
\begin{equation}
\text{tick}: \quad \underbrace{\text{phase observation}}_{\text{O}} \;\longrightarrow\; \underbrace{\text{categorical state count}}_{\text{C}} \;\longrightarrow\; \underbrace{\text{phase advance}}_{\text{P}}
\label{eq:tick}
\end{equation}
The three sub-operations are \emph{simultaneous}: the oscillator completes one period, which constitutes an observation of the partition cell's state, increments the count (computation), and advances the oscillator phase (physical processing). The arrow notation indicates logical decomposition, not temporal sequence.
\end{definition}

\begin{theorem}[Fundamental Identity of the Counting Loop]
\label{thm:counting-identity}
For any counting loop $\mathcal{L}$, a single tick simultaneously performs observation, computation, and processing. Formally, let $O(\mathcal{L})$ denote observation, $C(\mathcal{L})$ denote computation, and $P(\mathcal{L})$ denote physical processing. Then
\begin{equation}
O(\mathcal{L}) \equiv C(\mathcal{L}) \equiv P(\mathcal{L}).
\end{equation}
\end{theorem}

\begin{proof}
We prove the two identities separately.

\emph{Identity 1: $O(\mathcal{L}) \equiv C(\mathcal{L})$.} The oscillator's period defines a time interval $\Delta t = 2\pi/\omega$. During this interval, the oscillator's phase trajectory samples the partition cell $\mathcal{P}$. The number of ticks accumulated is the count $M$. By the ergodic theorem applied to the bounded oscillator~\cite{arnold1978}, the time average of any observable $A$ over $M$ ticks converges to the phase-space average:
\[
\langle A \rangle_{\text{time}} = \frac{1}{M} \sum_{k=1}^{M} A(t_k) \to \langle A \rangle_{\text{phase}}.
\]
The count $M$ therefore \emph{is} the observation---it encodes the state of $\mathcal{P}$ at resolution $M$. But $M$ is also the result of computation (counting). Hence observation and computation are identical operations: both produce $M$, and $M$ is the complete information content of both.

\emph{Identity 2: $C(\mathcal{L}) \equiv P(\mathcal{L})$.} Each tick increments the count $M \to M+1$, which changes the categorical address from $(n, \ell, m, s)$ to $(n', \ell', m', s')$ according to the selection rules of the partition structure. But each tick also advances the oscillator's physical phase by $2\pi$, which changes its position in phase space. Since the categorical address \emph{is} the phase space address (by \Cref{def:partition-coords}), the computational state change and the physical state change are the same operation. The count increment $M \to M+1$ and the phase advance $\theta \to \theta + 2\pi$ are two descriptions of one event.
\end{proof}

\begin{theorem}[Counting Loop Rate Equation]
\label{thm:rate-equation}
For a counting loop with count $M$, oscillator frequency $\omega$, and mean partition crossing time $\langle \taup \rangle$, the rate equation
\begin{equation}
\frac{dM}{dt} = \frac{\omega}{2\pi / M} = \frac{1}{\langle \taup \rangle}
\label{eq:rate-equation}
\end{equation}
holds identically.
\end{theorem}

\begin{proof}
The left equality: the oscillator completes $\omega/(2\pi)$ cycles per second. At count $M$, the partition cell at depth $n$ has $C(n) = 2n^2$ subdivisions. The effective tick rate at count $M$ is $\omega M/(2\pi)$, and the time per tick is $2\pi/(M\omega)$. Therefore $dM/dt = \omega/(2\pi/M)$.

The right equality: the mean partition crossing time $\langle \taup \rangle$ is the average time for the system to traverse one partition cell. Since each tick corresponds to one crossing, we have $\langle \taup \rangle = dt/dM$, or equivalently $dM/dt = 1/\langle \taup \rangle$.

The three-way equality $dM/dt = \omega/(2\pi/M) = 1/\langle \taup \rangle$ is the \emph{rate equation form} of the fundamental identity: the counting rate (computation), the oscillator frequency ratio (observation), and the inverse crossing time (physical processing) are all the same quantity.
\end{proof}

\begin{corollary}[No Other Primitive Needed]
\label{cor:no-other-primitive}
Since the counting loop performs observation, computation, and processing simultaneously (\Cref{thm:counting-identity}), and since these are the only three operations required for any computing system~\cite{turing1936}, no other computational primitive is needed. The counting loop is \emph{computationally complete} for the trajectory completion architecture.
\end{corollary}

\subsection{Coupled Counting Loop Networks}

\begin{definition}[Coupled Counting Loop Network]
\label{def:coupled-network}
A \emph{coupled counting loop network} is a graph $\mathcal{G} = (\mathcal{V}, \mathcal{E})$ where each vertex $v \in \mathcal{V}$ is a counting loop $\mathcal{L}_v = (\omega_v, \mathcal{P}_v, \tau_v)$ and each edge $(u,v) \in \mathcal{E}$ represents a \emph{coupling} with strength $\kappa_{uv} \geq 0$ that modifies the oscillator frequencies according to:
\begin{equation}
\dot{\theta}_v = \omega_v + \sum_{u : (u,v) \in \mathcal{E}} \kappa_{uv} \sin(\theta_u - \theta_v)
\label{eq:kuramoto-coupling}
\end{equation}
where $\theta_v$ is the phase of oscillator $v$~\cite{kuramoto1984}.
\end{definition}

\begin{theorem}[Harmonic Coincidence Computation]
\label{thm:harmonic-computation}
A coupled counting loop network computes via \emph{harmonic coincidence}: the output of the computation is encoded in the phase-locked states of the oscillator network. Information transfer between counting loops occurs through frequency synchronization, not through data transfer over communication channels.
\end{theorem}

\begin{proof}
By Kuramoto's synchronization theory~\cite{kuramoto1984,strogatz2003}, the coupled oscillator system~\eqref{eq:kuramoto-coupling} exhibits a phase transition at critical coupling strength $\kappa_c = 2/(\pi g(0))$, where $g$ is the distribution of natural frequencies. Below $\kappa_c$, oscillators are incoherent (independent). Above $\kappa_c$, a macroscopic fraction phase-locks with rational frequency ratios $\omega_u / \omega_v = p/q$ for $p, q \in \mathbb{Z}$.

These frequency ratios encode \emph{harmonic relationships} between partition cells. A ratio $p:q$ means oscillator $u$ completes $p$ ticks for every $q$ ticks of oscillator $v$. This is a \emph{computation}: it establishes a definite numerical relationship between the two partition cells without any data being ``sent'' from one oscillator to another.

By \Cref{thm:counting-identity}, each oscillator's frequency shift due to coupling is simultaneously: (a) an observation of the other oscillators' phases (via the sine coupling term), (b) a computation of the phase relationship (the resulting frequency ratio), and (c) a physical realization of the coupled state (the oscillator's actual motion). Therefore the network computes via harmonic coincidence---simultaneous phase-locking---not via sequential data transfer.
\end{proof}

\begin{figure}[t]
\centering
\fbox{\parbox{0.9\columnwidth}{\centering\vspace{2.5cm}\textbf{Figure 1: The Counting Loop Primitive}\\\vspace{0.3cm}Diagram showing oscillator $\omega$, partition cell $\mathcal{P}$ with coordinates $(n,\ell,m,s)$, and tick map $\tau$. Arrows illustrate the simultaneous observation--computation--processing cycle. Phase trajectory shown as spiral in bounded phase space. Inset: the fundamental identity $O \equiv C \equiv P$ expressed as three views of the same tick event.\vspace{2.5cm}}}
\caption{The counting loop primitive. A single tick simultaneously (a) observes the partition cell's state via phase sampling, (b) computes the categorical count $M \to M+1$, and (c) advances the oscillator's physical phase by $2\pi$. The three operations are ontologically identical, not merely correlated. The oscillator need not be electronic---it can be mechanical (engine piston), electromagnetic (camera frame clock), or any periodic physical process.}
\label{fig:counting-loop}
\end{figure}


% ===================================================================
\section{Architecture Overview}
\label{sec:architecture}

\subsection{The Vehicle as Coupled Counting Loop Network}

\begin{theorem}[Identity Unification]
\label{thm:identity-unification}
In the trajectory completion architecture, there is no separation between ``computer'' and ``car.'' Formally, let $\mathcal{A}$ be the set of actuators (steering, engine, brakes), $\mathcal{S}$ the set of sensors (cameras, LiDAR, IMU), and $\mathcal{C}$ the set of processors (MCUs, GPUs). Then $\mathcal{A}$, $\mathcal{S}$, and $\mathcal{C}$ are all subsets of a single coupled counting loop network $\mathcal{G}$:
\begin{equation}
\mathcal{A} \cup \mathcal{S} \cup \mathcal{C} \subseteq \mathcal{V}(\mathcal{G}).
\end{equation}
Moreover, the partition cells of actuator, sensor, and processor counting loops overlap in the shared partition space $\Gamma$.
\end{theorem}

\begin{proof}
Every physical component of a vehicle is an oscillator or contains oscillatory elements:
\begin{itemize}[leftmargin=*]
\item \emph{Engine}: pistons oscillate at frequency $f_{\mathrm{eng}} \approx 10$--$100$ Hz (600--6000 RPM). The crankshaft position sensor already counts these oscillations.
\item \emph{Wheels}: rotate at $f_{\mathrm{wheel}} \approx 1$--$20$ Hz ($\sim 60$--$1200$ RPM). ABS wheel speed sensors count every tooth passing.
\item \emph{Steering}: the steering column is a torsional oscillator with resonant frequency $f_{\mathrm{steer}} \approx 1$--$5$ Hz. Electric power steering motors add a driven oscillation.
\item \emph{Brakes}: brake discs rotate with wheel frequency; brake caliper pistons oscillate during ABS modulation at $f_{\mathrm{ABS}} \approx 10$--$15$ Hz.
\item \emph{Cameras}: image sensors sample at $f_{\mathrm{cam}} \approx 30$--$60$ Hz, with pixel readout clocks at $\sim 10^7$ Hz.
\item \emph{LiDAR}: the rotating mirror/motor oscillates at $f_{\mathrm{LiDAR}} \approx 10$--$20$ Hz.
\item \emph{IMU}: MEMS accelerometers are micro-oscillators at $f_{\mathrm{IMU}} \sim 10^3$--$10^4$ Hz.
\item \emph{Processors}: clock crystals oscillate at $f_{\mathrm{CPU}} \approx 10^8$--$10^{10}$ Hz.
\end{itemize}
Each oscillator defines a counting loop (\Cref{def:counting-loop}) by pairing it with the partition cell it occupies and the natural tick map. The couplings are physical: the engine drives the wheels through the drivetrain (mechanical coupling), the steering modifies wheel angles (kinematic coupling via tie rods), sensors sample the vehicle's state and environment (electromagnetic coupling), and processors read sensor data (electronic coupling via ADCs and buses). All these couplings have the Kuramoto form~\eqref{eq:kuramoto-coupling} to leading order (the coupling torque is proportional to a function of phase difference). Therefore the entire vehicle---metal, rubber, glass, silicon, and software---is a single coupled counting loop network $\mathcal{G}$, and the partition cells of all components lie in the shared phase space $\Gamma$.
\end{proof}

\subsection{Five Subsystems}

The architecture comprises five subsystems, adapted from the Buhera operating system~\cite{sachikonye2026buhera}:

\begin{enumerate}[leftmargin=*]
\item \textbf{Categorical State Manager (CSM)}: Maintains the current S-entropy triple $\Scurrent \in [0,1]^3$ from all Observer counting loops. Aggregates, weights, and fuses multi-perspective state estimates. Analog of Buhera's Categorical Memory Manager (CMM).

\item \textbf{Penultimate Navigation Engine (PNE)}: Computes the penultimate state $\Spen$ via backward trajectory completion through the ternary partition tree. Returns the unique cell one morphism from the destination. Analog of Buhera's Partition Space Scheduler (PSS).

\item \textbf{Sufficiency Recognition Module (SRM)}: Evaluates whether the transition $\Scurrent \to \Snext$ is safe via triple convergence of oscillatory, categorical, and partition perspectives. Returns \textsc{Proceed}, \textsc{Slow}, or \textsc{Stop}. Analog of Buhera's Dynamic Integrity Controller (DIC).

\item \textbf{Completion Morphism Executor (CME)}: Realizes the single-step transition $\Spen \to \Sfinal$ through the coupled oscillator network. The CME is not a controller---it is the physical coupling topology itself. Analog of Buhera's Partition Verification Engine (PVE).

\item \textbf{Triple Equivalence Monitor (TEM)}: Continuously verifies the rate equation $dM/dt = \omega/(2\pi/M) = 1/\langle \taup \rangle$ for all counting loops in the network. Detects violations within one tick period. Analog of Buhera's Triple Equivalence Monitor (TEM).
\end{enumerate}

\begin{figure}[t]
\centering
\fbox{\parbox{0.9\columnwidth}{\centering\vspace{3cm}\textbf{Figure 2: Five-Subsystem Architecture}\\\vspace{0.3cm}Block diagram showing CSM receiving inputs from Observer counting loops (wheel sensors, cameras, LiDAR, IMU, engine sensors). CSM feeds current state to PNE and SRM. PNE computes the penultimate state via backward navigation. SRM evaluates sufficiency through triple convergence. CME executes the completion morphism when SRM approves. TEM monitors all subsystems via rate equation verification. Coupling arrows represent oscillator couplings in the shared partition space, not data buses.\vspace{3cm}}}
\caption{The five-subsystem architecture. Unlike the conventional pipeline~\eqref{eq:conventional-pipeline}, information flows through oscillator couplings in the shared partition space, not through a sequential data pipeline. All subsystems operate simultaneously as coupled counting loop sub-networks within the vehicle's unified oscillator network.}
\label{fig:architecture}
\end{figure}

The formal interface contracts for each subsystem are specified in the following sections. Each contract defines preconditions, postconditions, and invariants in the style of design-by-contract~\cite{klein2009}, ensuring compositional correctness.


% ===================================================================
\section{Categorical State Manager}
\label{sec:csm}

The Categorical State Manager (CSM) is responsible for maintaining the current S-entropy triple $\Scurrent = (\Sk, \St, \Se)$ by aggregating counts from all Observer counting loops.

\begin{definition}[Observer Counting Loop]
\label{def:observer}
An \emph{Observer counting loop} $\mathcal{O}_i = (\omega_i, \mathcal{P}_i, \tau_i)$ is a counting loop whose partition cell $\mathcal{P}_i$ corresponds to a specific observable aspect of the vehicle's state or environment. The Observer produces a local S-entropy estimate $\Sentropy_i = (\Sk^{(i)}, \St^{(i)}, \Se^{(i)})$ derived from its tick count $M_i$.
\end{definition}

\begin{definition}[CSM State Estimation]
\label{def:csm-estimation}
Given $K$ Observer counting loops $\{\mathcal{O}_1, \ldots, \mathcal{O}_K\}$ with local S-entropy estimates $\{\Sentropy_1, \ldots, \Sentropy_K\}$ and reliability weights $\{w_1, \ldots, w_K\}$ satisfying $\sum_i w_i = 1$, the CSM computes the global state estimate:
\begin{equation}
\Scurrent = \sum_{i=1}^{K} w_i \, \Sentropy_i.
\label{eq:csm-estimate}
\end{equation}
The weights are inversely proportional to the variance of each Observer's recent estimates:
\begin{equation}
w_i = \frac{\sigma_i^{-2}}{\sum_{j=1}^{K} \sigma_j^{-2}}
\end{equation}
where $\sigma_i^2 = \mathrm{Var}[\Sentropy_i]$ is estimated from the last $W$ ticks of Observer $i$.
\end{definition}

This is the \emph{inverse-variance weighting} scheme, which is the minimum-variance unbiased estimator when the individual estimates are unbiased and uncorrelated~\cite{jaynes1957}. The scheme naturally downweights noisy sensors and upweights reliable ones, with no hand-tuning required.

\begin{theorem}[CSM Convergence]
\label{thm:csm-convergence}
With $K$ independent Observer counting loops, the CSM estimation error satisfies
\begin{equation}
\|\Scurrent - \Sentropy_{\mathrm{true}}\| \leq \frac{C}{\sqrt{K}}
\end{equation}
for a constant $C$ depending on the maximum individual Observer variance.
\end{theorem}

\begin{proof}
Since the weights are chosen to minimize variance (inverse-variance weighting is optimal in the Gauss--Markov sense), the combined estimator has variance
\[
\mathrm{Var}[\Scurrent] = \left(\sum_{i=1}^K \sigma_i^{-2}\right)^{-1} \leq \frac{\sigma_{\max}^2}{K}
\]
where $\sigma_{\max}^2 = \max_i \sigma_i^2$. The estimation error in any coordinate is
\[
|S_{\mathrm{current},j} - S_{\mathrm{true},j}| \leq \frac{\sigma_{\max}}{\sqrt{K}} \cdot z_\alpha
\]
with probability $1-\alpha$, where $z_\alpha$ is the standard normal quantile. Taking the $\ell^2$ norm across the three S-entropy coordinates and setting $C = \sqrt{3} \cdot \sigma_{\max} \cdot z_\alpha$ gives $\|\Scurrent - \Sentropy_{\mathrm{true}}\| \leq C/\sqrt{K}$.
\end{proof}

\begin{definition}[Urgency Tiers]
\label{def:urgency}
Let $M$ be the current partition depth and $\dcat(\Scurrent, \Snext)$ the categorical distance between the current and next states. The CSM assigns urgency tiers:
\begin{align}
\text{Tier 1 (immediate):} &\quad \dcat < 3^{-M+1} \label{eq:tier1}\\
\text{Tier 2 (approaching):} &\quad 3^{-M+1} \leq \dcat < 3^{-M+3} \label{eq:tier2}\\
\text{Tier 3 (strategic):} &\quad \dcat \geq 3^{-M+3} \label{eq:tier3}
\end{align}
\end{definition}

The tier thresholds are ternary powers because the partition tree has branching factor~3. Tier~1 means the next state is within the same parent cell (one level up the tree): e.g., a lane-keeping adjustment within the current road segment. Tier~2 means it is within the grandparent's subtree: e.g., a lane change or turn within the current neighborhood. Tier~3 means it requires navigating higher in the tree: e.g., a highway exit or route change.

The urgency tier determines the SRM's tolerance $\delta$: Tier~1 transitions have the tightest tolerance (highest confidence required for safety), Tier~3 transitions have the loosest (more time to correct). This mirrors human driving: immediate maneuvers require high confidence, while strategic decisions tolerate more uncertainty.

\textbf{Formal Interface Contract.}
\begin{itemize}[leftmargin=*]
\item \emph{Precondition}: At least one Observer counting loop is operational ($K \geq 1$).
\item \emph{Postcondition}: $\Scurrent \in [0,1]^3$ is updated every $\Delta t_{\mathrm{CSM}} \leq 1/f_{\min}$ seconds, where $f_{\min}$ is the minimum Observer frequency. Each component satisfies $0 \leq S_j \leq 1$.
\item \emph{Invariant}: The state change rate is bounded: $\|\Scurrent(t+\Delta t) - \Scurrent(t)\| \leq v_{\max} \Delta t / L$ where $v_{\max}$ is maximum vehicle speed and $L$ is the partition cell diameter at current depth. Violation of this bound triggers TEM alert.
\end{itemize}


% ===================================================================
\section{Penultimate Navigation Engine}
\label{sec:pne}

The Penultimate Navigation Engine (PNE) replaces conventional path planning (A*, Dijkstra, RRT*, D*-Lite) with \emph{backward navigation} through the ternary partition tree.

\subsection{Road Network as Partition Tree}

\begin{definition}[Road Network Partition Tree]
\label{def:road-tree}
The road network is represented as a ternary partition tree $\mathcal{T}$ of depth $D$:
\begin{itemize}[leftmargin=*]
\item Depth 0: entire road network (root node).
\item Depths 1--3: highway-level segments (interstate, Autobahn, motorway).
\item Depths 4--6: arterial-level segments (main roads, boulevards).
\item Depths 7--9: local street segments (residential, Stra{\ss}en).
\item Depths $\geq 10$: lane-level and sub-lane resolution (individual lanes, positions within lane).
\end{itemize}
Each node at depth $d$ has at most 3 children at depth $d+1$. The total number of leaf nodes is $N \leq 3^D$.
\end{definition}

The ternary branching factor is not arbitrary: it corresponds to the categorical trichotomy of each decision (left/center/right at each resolution level). A highway segment splits into three arterial regions; an arterial splits into three local regions; a local street splits into three lane-level positions.

\begin{definition}[Penultimate State]
\label{def:penultimate}
Given a destination state $\Sfinal$ (a leaf in the partition tree $\mathcal{T}$), the \emph{penultimate state} $\Spen$ is the unique partition cell that is exactly one morphism (one tree edge) away from $\Sfinal$ on the path from $\Scurrent$:
\begin{equation}
\Spen = \mathrm{path}(\Scurrent, \Sfinal)[-2]
\end{equation}
where $\mathrm{path}(\Scurrent, \Sfinal)$ is the unique tree path and $[-2]$ denotes the second-to-last element.
\end{definition}

\subsection{The Backward Navigation Algorithm}

\begin{algorithm}[t]
\caption{BackwardNavigation($\Sfinal$, $\mathcal{T}$, $\Scurrent$)}
\label{alg:backward-nav}
\begin{algorithmic}[1]
\REQUIRE Destination $\Sfinal$, partition tree $\mathcal{T}$, current state $\Scurrent$
\ENSURE Penultimate state $\Spen$
\STATE $v_{\mathrm{dest}} \leftarrow \mathrm{leaf}(\Sfinal, \mathcal{T})$ \COMMENT{Locate destination leaf: $\BigO(D)$}
\STATE $v_{\mathrm{curr}} \leftarrow \mathrm{leaf}(\Scurrent, \mathcal{T})$ \COMMENT{Locate current leaf: $\BigO(D)$}
\STATE $v_{\mathrm{lca}} \leftarrow \mathrm{LCA}(v_{\mathrm{dest}}, v_{\mathrm{curr}}, \mathcal{T})$ \COMMENT{Lowest common ancestor: $\BigO(D)$}
\STATE $\mathrm{path}_{\mathrm{up}} \leftarrow$ trace from $v_{\mathrm{curr}}$ to $v_{\mathrm{lca}}$
\STATE $\mathrm{path}_{\mathrm{down}} \leftarrow$ trace from $v_{\mathrm{lca}}$ to $v_{\mathrm{dest}}$
\STATE $\mathrm{path} \leftarrow \mathrm{path}_{\mathrm{up}} \circ \mathrm{path}_{\mathrm{down}}$ \COMMENT{Concatenate}
\STATE $\Spen \leftarrow \mathrm{path}[\mathrm{len}(\mathrm{path}) - 2]$ \COMMENT{Penultimate node}
\RETURN $\Spen$
\end{algorithmic}
\end{algorithm}

\begin{theorem}[Navigation Complexity]
\label{thm:nav-complexity}
\Cref{alg:backward-nav} computes the penultimate state in time $\BigO(\log_3 N)$ where $N$ is the number of road network leaf segments.
\end{theorem}

\begin{proof}
The tree $\mathcal{T}$ has depth $D = \lceil\log_3 N\rceil$. Each of the operations in lines 1--7 requires at most $2D$ steps:
\begin{itemize}[leftmargin=*]
\item Leaf location (lines 1--2): traverse from root to leaf by comparing coordinates at each depth level. Each comparison selects one of three children: $\BigO(D)$ comparisons.
\item LCA computation (line 3): trace both leaves upward until their paths meet. At most $D$ steps per leaf, total $\BigO(2D) = \BigO(D)$.
\item Path construction (lines 4--6): record the $\BigO(D)$ nodes on each half-path and concatenate: $\BigO(D)$.
\item Penultimate extraction (line 7): array indexing: $\BigO(1)$.
\end{itemize}
Total: $\BigO(D) = \BigO(\log_3 N)$.
\end{proof}

\begin{theorem}[Correctness of Backward Navigation]
\label{thm:nav-correctness}
\Cref{alg:backward-nav} returns the unique penultimate state: the partition cell that is exactly one completion morphism from the destination along the tree path from the current position.
\end{theorem}

\begin{proof}
In the ternary partition tree $\mathcal{T}$, there is a unique path between any two nodes (since a tree is connected and acyclic). The path from $v_{\mathrm{curr}}$ to $v_{\mathrm{dest}}$ passes through their lowest common ancestor $v_{\mathrm{lca}}$ and is unique. The penultimate state is the second-to-last node on this unique path, which is therefore itself unique. By construction (line 7), the algorithm returns exactly this node. Since the last edge on the path connects $\Spen$ to $\Sfinal$, the morphism $\varphi: \Spen \to \Sfinal$ is a single tree edge---one completion morphism.
\end{proof}

\begin{proposition}[Dynamic Re-Navigation]
\label{prop:dynamic-renav}
If the road network changes (e.g., lane closure, accident, construction), re-navigation requires at most $\BigO(\log_3 N)$ time.
\end{proposition}

\begin{proof}
A road network change modifies $\BigO(1)$ nodes in $\mathcal{T}$ (the affected segments and their ancestors up to the root, but ancestors are $\BigO(D)$ nodes). The LCA may change, but the entire backward navigation algorithm is re-executed from scratch in $\BigO(D) = \BigO(\log_3 N)$ time. No incremental graph update is needed because the algorithm's complexity is already logarithmic.
\end{proof}

\textbf{Comparison with conventional path planning.} Consider a city road network with $V = 10^5$ intersections and $E = 3 \times 10^5$ road segments:
\begin{itemize}[leftmargin=*]
\item A* / Dijkstra: $\BigO(E + V \log V) \approx 10^6$ operations.
\item RRT*: $\BigO(n \log n)$ where $n$ is sample count, typically $n \sim 10^4$: $\approx 10^5$ operations.
\item Backward navigation with $N = 10^5$ leaves: $\BigO(\log_3 10^5) = \BigO(10.5) \approx 11$ operations.
\end{itemize}
The improvement is not incremental---it is five orders of magnitude. This is possible because backward navigation exploits the hierarchical partition structure, while graph-search algorithms treat the network as an unstructured graph.

\begin{figure}[t]
\centering
\fbox{\parbox{0.9\columnwidth}{\centering\vspace{3cm}\textbf{Figure 3: Backward Navigation in the Ternary Partition Tree}\\\vspace{0.3cm}Illustration of a ternary partition tree with depth $D=9$ representing a city road network. The tree is drawn with root at top. Current state $\Scurrent$ (left leaf, highlighted blue) and destination $\Sfinal$ (right leaf, highlighted green) are connected by a path through their lowest common ancestor $v_{\mathrm{lca}}$ (highlighted yellow). The penultimate state $\Spen$ is circled in red, one edge before $\Sfinal$. Depth levels are labeled: highway (1--3), arterial (4--6), local (7--9). Faded subtrees show pruned branches not on the navigation path.\vspace{3cm}}}
\caption{Backward navigation in the ternary partition tree. The unique path from $\Scurrent$ to $\Sfinal$ passes through the lowest common ancestor. The penultimate state $\Spen$ is the unique cell one completion morphism from the destination. The algorithm visits only $\BigO(\log_3 N)$ nodes (bold path), ignoring the vast majority of the tree.}
\label{fig:backward-nav}
\end{figure}

\textbf{Formal Interface Contract.}
\begin{itemize}[leftmargin=*]
\item \emph{Precondition}: $\Sfinal$ is a valid leaf in $\mathcal{T}$; $\Scurrent$ is provided by CSM with convergence guarantee (\Cref{thm:csm-convergence}).
\item \emph{Postcondition}: $\Spen$ is the unique penultimate state; $\dcat(\Spen, \Sfinal) = 3^{-D+1}$ (one leaf-level separation in the tree metric).
\item \emph{Invariant}: If $\Scurrent$ changes (vehicle moves), $\Spen$ is recomputed within $\BigO(\log_3 N)$ ticks of the fastest Observer. The penultimate state is always current.
\end{itemize}


% ===================================================================
\section{Sufficiency Recognition Module}
\label{sec:srm}

The Sufficiency Recognition Module (SRM) is the most critical subsystem of the architecture. It replaces \emph{prediction} with \emph{recognition}: rather than forecasting what will happen and planning around the forecast, it assesses whether the current state affords sufficient margin for the next transition. This is the formal counterpart of what human drivers do intuitively.

\subsection{Information Sufficiency}

\begin{definition}[Thermodynamic Gap]
\label{def:thermo-gap}
The \emph{thermodynamic gap} for a proposed transition $\Scurrent \to \Snext$ is the free energy difference:
\begin{equation}
\varepsilon = F(\Scurrent) - F(\Snext) - W_{\mathrm{min}}
\end{equation}
where $F(\Sentropy) = E(\Sentropy) - T \cdot H(\Sentropy)$ is the Helmholtz free energy of the state, $E$ is the internal energy, $T$ is the effective temperature, $H$ is the entropy, and $W_{\mathrm{min}}$ is the minimum work required for the transition~\cite{landauer1961,bennett1982,jaynes1957}.
\end{definition}

A positive gap $\varepsilon > 0$ means the transition is thermodynamically favorable: more free energy is available than the transition requires. A negative gap $\varepsilon < 0$ means the transition requires external energy input beyond what is available. A zero gap $\varepsilon = 0$ is the critical boundary.

\begin{definition}[Three-Perspective Measurement]
\label{def:three-perspective}
The SRM measures the thermodynamic gap from three independent perspectives, each corresponding to one term of the fundamental identity $O \equiv C \equiv P$:
\begin{align}
\varepsilon_{\mathrm{osc}} &= \text{gap from oscillatory (frequency/observation) perspective} \label{eq:eps-osc} \\
\varepsilon_{\mathrm{cat}} &= \text{gap from categorical (count/computation) perspective} \label{eq:eps-cat} \\
\varepsilon_{\mathrm{par}} &= \text{gap from partition (geometry/processing) perspective} \label{eq:eps-par}
\end{align}
Specifically, the three gap measurements are:
\begin{align}
\varepsilon_{\mathrm{osc}} &= \frac{\hbar}{2}\sum_{v \in \mathcal{V}} \left(\omega_v^{(\mathrm{next})} - \omega_v^{(\mathrm{current})}\right) - W_{\mathrm{min}}^{(\mathrm{osc})} \label{eq:eps-osc-formula}\\
\varepsilon_{\mathrm{cat}} &= k_B T \ln\!\left(\frac{C(n_{\mathrm{next}})}{C(n_{\mathrm{current}})}\right) - W_{\mathrm{min}}^{(\mathrm{cat})} \label{eq:eps-cat-formula}\\
\varepsilon_{\mathrm{par}} &= \frac{1}{2}k_B T \left(\|\Snext\|^2 - \|\Scurrent\|^2\right) - W_{\mathrm{min}}^{(\mathrm{par})} \label{eq:eps-par-formula}
\end{align}
where the three minimum-work terms $W_{\mathrm{min}}^{(\cdot)}$ account for the irreducible cost of the transition in each perspective, and $C(n) = 2n^2$ is the partition degeneracy~\eqref{eq:degeneracy}.
\end{definition}

The oscillatory perspective~\eqref{eq:eps-osc-formula} measures the energy change through frequency shifts of the vehicle's oscillators. The categorical perspective~\eqref{eq:eps-cat-formula} measures the entropy change through the ratio of available states at the current and next depths. The partition perspective~\eqref{eq:eps-par-formula} measures the geometric change through the S-entropy norm difference. By the fundamental identity, all three should agree; their agreement is what constitutes sufficiency recognition.

\begin{definition}[Triple Convergence]
\label{def:triple-convergence}
The three-perspective measurement satisfies \emph{triple convergence} with tolerance $\delta > 0$ if
\begin{equation}
|\varepsilon_{\mathrm{osc}} - \varepsilon_{\mathrm{cat}}| < \delta \quad \text{and} \quad |\varepsilon_{\mathrm{cat}} - \varepsilon_{\mathrm{par}}| < \delta.
\label{eq:triple-convergence}
\end{equation}
\end{definition}

Note that only two pairwise comparisons are needed: by the triangle inequality, $|\varepsilon_{\mathrm{osc}} - \varepsilon_{\mathrm{par}}| \leq |\varepsilon_{\mathrm{osc}} - \varepsilon_{\mathrm{cat}}| + |\varepsilon_{\mathrm{cat}} - \varepsilon_{\mathrm{par}}| < 2\delta$.

\begin{definition}[Information Sufficiency]
\label{def:info-sufficiency}
A transition $\Scurrent \to \Snext$ is \emph{informationally sufficient} if:
\begin{enumerate}[leftmargin=*]
\item Triple convergence holds (\Cref{def:triple-convergence}), and
\item The mean gap $\bar{\varepsilon} = (\varepsilon_{\mathrm{osc}} + \varepsilon_{\mathrm{cat}} + \varepsilon_{\mathrm{par}})/3$ satisfies $\bar{\varepsilon} > 0$.
\end{enumerate}
The first condition ensures the measurement is \emph{reliable} (three perspectives agree). The second ensures the transition is \emph{feasible} (positive thermodynamic driving force).
\end{definition}

\subsection{The Sufficiency Theorem}

\begin{theorem}[Sufficiency Theorem]
\label{thm:sufficiency}
If the transition $\Scurrent \to \Snext$ satisfies information sufficiency (\Cref{def:info-sufficiency})---triple convergence with tolerance $\delta$ and mean gap $\bar{\varepsilon} > 0$---then the transition is safe.
\end{theorem}

\begin{proof}
The proof proceeds in three steps.

\textbf{Step 1: Triple equivalence grounds the measurement.} By the fundamental identity $O(x) \equiv C(x) \equiv P(x)$ (\Cref{thm:counting-identity}), the oscillatory, categorical, and partition perspectives describe the same physical reality from three independent vantage points. The oscillatory perspective measures through temporal frequency, the categorical through combinatorial counting, and the partition through spatial geometry. They use different physical mechanisms (frequency measurement, state enumeration, distance computation) but refer to the same underlying thermodynamic gap. The three measurements $\varepsilon_{\mathrm{osc}}, \varepsilon_{\mathrm{cat}}, \varepsilon_{\mathrm{par}}$ are therefore three independent estimates of the same physical quantity $\varepsilon_{\mathrm{true}}$.

\textbf{Step 2: Convergence confirms recognition with quantifiable confidence.} By the G\"{o}delian residue theorem~\cite{godel1931,sachikonye2026trajectory}, any self-referential measurement has irreducible uncertainty $\varepsilon > 0$---the system cannot know itself perfectly from within. However, if three \emph{independent} measurements agree to within $\delta$, the probability of coincidental agreement (all three wrong in the same direction by the same amount, producing false convergence) is bounded.

Let $\varepsilon_{\mathrm{true}}$ be the true gap and suppose the transition is actually unsafe ($\varepsilon_{\mathrm{true}} < 0$). For the SRM to falsely declare it safe, all three measurements must be wrong by at least $|\varepsilon_{\mathrm{true}}|$, and must agree to within $\delta$. By \Cref{lem:observer-independence} (proved below), the measurements are independent. The probability that a single measurement deviates by more than $|\varepsilon_{\mathrm{true}}|$ and lands in a window of width $\delta$ is at most $\delta / |\varepsilon_{\mathrm{true}}|$ (uniform bound). For two independent pairwise constraints:
\begin{equation}
P_{\mathrm{false}} \leq \left(\frac{\delta}{|\varepsilon_{\mathrm{true}}|}\right)^2.
\end{equation}
With high-frequency counting loops ($\omega \sim 10^4$--$10^{10}$ Hz), the measurement resolution is $\delta \sim 10^{-6}|\varepsilon_{\mathrm{true}}|$, giving $P_{\mathrm{false}} \sim 10^{-12}$---far below any engineering safety requirement.

\textbf{Step 3: Positive gap ensures thermodynamic feasibility.} The condition $\bar{\varepsilon} > 0$ ensures the transition has a net thermodynamic driving force---the available free energy exceeds the required work. By the second law of thermodynamics~\cite{boltzmann1896}, spontaneous transitions proceed in the direction of decreasing free energy. The positive gap provides a safety margin: even in the worst case (true gap is $\bar{\varepsilon} - 2\delta$ given convergence tolerance), the transition remains favorable as long as $\bar{\varepsilon} > 2\delta$. The SRM enforces this by requiring $\bar{\varepsilon} > \varepsilon_{\mathrm{threshold}} \geq 2\delta$.
\end{proof}

\begin{theorem}[Non-Convergence Safety]
\label{thm:non-convergence}
If triple convergence fails (i.e., condition~\eqref{eq:triple-convergence} is violated), then stopping the vehicle is always a safe response.
\end{theorem}

\begin{proof}
A vehicle at rest has velocity $\mathbf{v} = 0$ and acceleration $\mathbf{a} = 0$. Its temporal entropy is $\St = 0$ (zero rate of state change). The S-entropy of the rest state is $\Sentropy_{\mathrm{rest}} = (\Sk, 0, \Se)$ for some fixed categorical and energetic entropies.

We prove the rest state is thermodynamically stable. By the second law, the total entropy of the vehicle-environment system satisfies $dS_{\mathrm{total}}/dt \geq 0$~\cite{boltzmann1896}. The vehicle at rest has minimum kinetic contribution to entropy. To leave the rest state requires a positive change $\Delta \St > 0$ (the vehicle must start moving), which requires work input $W \geq k_B T \ln(C(n_{\mathrm{next}})/C(n_{\mathrm{current}}))$ by the categorical perspective~\eqref{eq:eps-cat-formula}. Without the SRM confirming sufficiency (triple convergence has failed), the CME does not execute any morphism, so no work is directed into a transition. The vehicle therefore remains at rest.

The rest state is stable against perturbations: small perturbations (road vibration, wind) produce small $\St > 0$ but the braking system (a damping oscillator) returns $\St \to 0$. The braking system does not require SRM approval to maintain rest---it is a passive physical damper.

Therefore, when the SRM cannot confirm a transition is safe (non-convergence), commanding a stop is always safe.
\end{proof}

\begin{corollary}[Uncertainty Implies Deceleration]
\label{cor:uncertainty-deceleration}
If triple convergence holds but the mean gap $\bar{\varepsilon}$ is small ($0 < \bar{\varepsilon} < \varepsilon_{\mathrm{threshold}}$), the SRM commands reduced speed. This matches human driving behavior: uncertainty causes deceleration. Formally, the SRM maps $\bar{\varepsilon}$ to a maximum safe speed via $v_{\mathrm{safe}} \propto \sqrt{\bar{\varepsilon}}$ (kinetic energy scales as $v^2$).
\end{corollary}

\begin{algorithm}[t]
\caption{SufficiencyCheck($\Scurrent$, $\Snext$, $\{\mathcal{O}_i\}_{i=1}^K$)}
\label{alg:sufficiency}
\begin{algorithmic}[1]
\REQUIRE Current state $\Scurrent$, proposed next state $\Snext$, Observers $\{\mathcal{O}_i\}$
\ENSURE Decision $d \in \{\textsc{Proceed}, \textsc{Slow}, \textsc{Stop}\}$
\STATE $\varepsilon_{\mathrm{osc}} \leftarrow$ ComputeOscillatoryGap($\Scurrent$, $\Snext$, $\{\mathcal{O}_i\}_{\mathrm{osc}}$)
\STATE $\varepsilon_{\mathrm{cat}} \leftarrow$ ComputeCategoricalGap($\Scurrent$, $\Snext$, $\{\mathcal{O}_i\}_{\mathrm{cat}}$)
\STATE $\varepsilon_{\mathrm{par}} \leftarrow$ ComputePartitionGap($\Scurrent$, $\Snext$, $\{\mathcal{O}_i\}_{\mathrm{par}}$)
\STATE $\delta \leftarrow$ GetTolerance(Tier($\Scurrent$, $\Snext$)) \COMMENT{Tier-dependent}
\IF{$|\varepsilon_{\mathrm{osc}} - \varepsilon_{\mathrm{cat}}| \geq \delta$ \textbf{or} $|\varepsilon_{\mathrm{cat}} - \varepsilon_{\mathrm{par}}| \geq \delta$}
    \RETURN \textsc{Stop} \COMMENT{Non-convergence: safe halt}
\ENDIF
\STATE $\bar{\varepsilon} \leftarrow (\varepsilon_{\mathrm{osc}} + \varepsilon_{\mathrm{cat}} + \varepsilon_{\mathrm{par}})/3$
\IF{$\bar{\varepsilon} < 0$}
    \RETURN \textsc{Stop} \COMMENT{Negative gap: transition infeasible}
\ELSIF{$\bar{\varepsilon} < \varepsilon_{\mathrm{threshold}}$}
    \RETURN \textsc{Slow} \COMMENT{Low margin: reduce speed}
\ENDIF
\RETURN \textsc{Proceed} \COMMENT{Triple convergence confirmed, sufficient margin}
\end{algorithmic}
\end{algorithm}

\subsection{Formal Safety Analysis of the SRM}

\begin{lemma}[Observer Independence]
\label{lem:observer-independence}
The oscillatory, categorical, and partition perspectives are informationally independent: no single-point failure can corrupt all three measurements simultaneously.
\end{lemma}

\begin{proof}
The three perspectives measure fundamentally different physical quantities through different mechanisms:
\begin{itemize}[leftmargin=*]
\item $\varepsilon_{\mathrm{osc}}$ measures frequency differences via temporal correlations in oscillator ticks. It depends on the \emph{timing} of ticks.
\item $\varepsilon_{\mathrm{cat}}$ measures state count ratios via combinatorial enumeration. It depends on the \emph{number} of occupied cells.
\item $\varepsilon_{\mathrm{par}}$ measures geometric distances in S-entropy space. It depends on the \emph{coordinates} of the current and next states.
\end{itemize}
Consider potential failure modes:
\begin{enumerate}
\item \emph{Sensor failure} (e.g., camera blackout): affects the Observers contributing to one or two perspectives, but not all three (the wheel speed sensor contributes to $\varepsilon_{\mathrm{osc}}$ but not $\varepsilon_{\mathrm{par}}$; the IMU contributes to $\varepsilon_{\mathrm{par}}$ but not $\varepsilon_{\mathrm{cat}}$).
\item \emph{Computational error} (e.g., bit flip): affects the counting mechanism ($\varepsilon_{\mathrm{cat}}$) but not the analog frequency measurement ($\varepsilon_{\mathrm{osc}}$) or the physical geometry ($\varepsilon_{\mathrm{par}}$).
\item \emph{Physical perturbation} (e.g., pothole): affects $\varepsilon_{\mathrm{par}}$ (geometric change) but is \emph{correctly captured} by $\varepsilon_{\mathrm{osc}}$ (frequency shift in suspension oscillators) and $\varepsilon_{\mathrm{cat}}$ (count change in affected cells).
\end{enumerate}
In no case does a single failure mode corrupt all three perspectives in the same direction by the same amount. The perspectives are therefore informationally independent for the purpose of false-positive analysis.
\end{proof}

\begin{theorem}[SRM False-Positive Bound]
\label{thm:srm-false-positive}
The probability that the SRM declares a transition safe when it is actually unsafe is bounded by
\begin{equation}
P_{\mathrm{false\text{-}positive}} \leq \left(\frac{\delta}{\varepsilon_{\min}}\right)^2
\end{equation}
where $\varepsilon_{\min}$ is the minimum true gap magnitude for unsafe transitions and $\delta$ is the convergence tolerance.
\end{theorem}

\begin{proof}
By \Cref{lem:observer-independence}, the three gap measurements are informationally independent. For a false positive to occur, the transition must be unsafe ($\varepsilon_{\mathrm{true}} < -\varepsilon_{\min}$) while the three measurements all report $\varepsilon > 0$ and agree to within $\delta$.

Each measurement independently must deviate from the true value by at least $\varepsilon_{\min}$ (to change sign) and land in a window of width $\delta$ (to satisfy convergence). The probability of a single measurement deviating by $\varepsilon_{\min}$ and hitting a $\delta$-window is bounded by $\delta/\varepsilon_{\min}$ under any sub-Gaussian noise model. For the two pairwise convergence constraints~\eqref{eq:triple-convergence} to be simultaneously satisfied with all three measurements wrong, the joint probability is at most $(\delta/\varepsilon_{\min})^2$ by independence.
\end{proof}

\textbf{Formal Interface Contract.}
\begin{itemize}[leftmargin=*]
\item \emph{Precondition}: $\Scurrent$ from CSM with $\|\Scurrent - \Sentropy_{\mathrm{true}}\| \leq C/\sqrt{K}$; $\Snext$ from PNE; at least one Observer contributing to each of the three perspectives.
\item \emph{Postcondition}: Decision $d \in \{\textsc{Proceed}, \textsc{Slow}, \textsc{Stop}\}$ with $P_{\mathrm{false\text{-}positive}} \leq (\delta/\varepsilon_{\min})^2$.
\item \emph{Invariant}: Non-convergence \emph{always} maps to \textsc{Stop}; never to \textsc{Proceed} or \textsc{Slow}. This invariant is unconditional---it holds regardless of any other system state.
\end{itemize}


% ===================================================================
\section{Completion Morphism Executor}
\label{sec:cme}

The Completion Morphism Executor (CME) is not a controller in the conventional control-theoretic sense. It does not compute a control signal from state error and apply it to actuators. Instead, it \emph{is} the coupled oscillator network itself, executing the single-step transition from $\Spen$ to $\Sfinal$ through a phase transition in the oscillator network.

\subsection{The Completion Morphism as Phase Transition}

\begin{theorem}[Completion Morphism as Phase Transition]
\label{thm:morphism-phase-transition}
A completion morphism $\varphi: \Spen \to \Sfinal$ corresponds to a phase transition in the coupled oscillator network: the system transitions from one phase-locked state (encoding $\Spen$) to another (encoding $\Sfinal$) through a transient desynchronization.
\end{theorem}

\begin{proof}
In the phase-locked state corresponding to $\Spen$, the oscillators have frequency ratios $\{\omega_i/\omega_j\}_{(i,j) \in \mathcal{E}}$ that encode the penultimate partition cell's categorical coordinates. These ratios are rational numbers determined by the partition structure: $\omega_i/\omega_j = p_{ij}/q_{ij}$ where $p_{ij}, q_{ij} \in \mathbb{Z}$ correspond to the harmonic relationship between partition cells $\mathcal{P}_i$ and $\mathcal{P}_j$.

The completion morphism $\varphi$ changes the target coordinates from $\Spen$ to $\Sfinal$, which changes the target frequency ratios. By Kuramoto synchronization theory~\cite{kuramoto1984}, when the target ratios change, the coupled system undergoes:
\begin{enumerate}
\item \emph{Desynchronization}: the current phase-locked state becomes unstable as the coupling landscape shifts.
\item \emph{Transient}: oscillators adjust their phases through the coupling forces $\kappa_{uv}\sin(\theta_u - \theta_v)$.
\item \emph{Resynchronization}: the system relaxes to the new phase-locked state encoding $\Sfinal$.
\end{enumerate}
This sequence is precisely a phase transition: the order parameter (Kuramoto coherence $r = |N^{-1}\sum_j e^{i\theta_j}|$) drops during desynchronization and recovers during resynchronization~\cite{strogatz2003}. The timescale is determined by the coupling strengths $\kappa_{ij}$ and the frequency differences $|\omega_i - \omega_j|$, not by any explicit trajectory computation.
\end{proof}

\begin{definition}[Morphism Types]
\label{def:morphism-types}
Completion morphisms are classified by the dominant S-entropy coordinate change:
\begin{itemize}[leftmargin=*]
\item \textbf{Longitudinal morphism ($\Delta \St$ dominant)}: The primary change is in temporal entropy. Engine and brake oscillators shift their phase relationship. Concretely: the engine frequency increases (acceleration via throttle opening, changing engine RPM) or the brake applies friction torque (deceleration, introducing a damping oscillation). The drivetrain coupling $\kappa_{\mathrm{drive}}$ transmits the phase change from engine to wheels.

\item \textbf{Lateral morphism ($\Delta \Sk$ dominant)}: The primary change is in categorical entropy. The steering oscillator shifts its phase relationship with the wheel oscillators. Concretely: the steering motor applies a torsional force, changing the wheel angle. The tie-rod coupling $\kappa_{\mathrm{steer}}$ transmits the phase change from steering column to wheel assemblies. The vehicle moves between partition cells in the lateral direction (lane change, turn).

\item \textbf{Energy morphism ($\Delta \Se$ dominant)}: The primary change is in energetic entropy. The transmission oscillator changes its frequency ratio (gear shift). Concretely: the transmission changes the gear ratio $r_g$, which changes the coupling ratio between engine and wheel oscillators from $\omega_{\mathrm{eng}}/\omega_{\mathrm{wheel}} = r_{g,\mathrm{old}}$ to $r_{g,\mathrm{new}}$. The energy distribution between kinetic and potential forms changes.
\end{itemize}
\end{definition}

In practice, most driving morphisms are combinations. A highway merge is primarily longitudinal (accelerate) with a lateral component (lane change). A hill climb is primarily energy (gear downshift) with a longitudinal component (speed adjustment). The S-entropy triple captures these combinations naturally as vectors in $[0,1]^3$.

\begin{theorem}[Coupling Topology Determines Control Law]
\label{thm:coupling-control}
In the trajectory completion architecture, the physical coupling topology between oscillators \emph{is} the control law. No separate control algorithm exists or is needed.
\end{theorem}

\begin{proof}
By the fundamental identity $O(x) \equiv C(x) \equiv P(x)$ (\Cref{thm:counting-identity}), the physical coupling between oscillators simultaneously constitutes:
\begin{enumerate}
\item \emph{Observation}: the coupling force $\kappa_{uv}\sin(\theta_u - \theta_v)$ transmits phase information from oscillator $u$ to oscillator $v$. This is how each oscillator ``observes'' the others.
\item \emph{Computation}: the resulting phase relationships $\{\theta_u - \theta_v\}$ encode the vehicle's state in partition coordinates. The computation of ``where am I and where should I be'' is performed by the network's relaxation to equilibrium.
\item \emph{Processing}: the coupling forces physically realize the state transition. The torque applied to the steering column by the coupling is the ``control signal.''
\end{enumerate}

A conventional PID controller computes $u(t) = K_p e(t) + K_i \int e \, dt + K_d \dot{e}(t)$ where $e = x_{\mathrm{ref}} - x$. In the coupled oscillator framework:
\begin{itemize}[leftmargin=*]
\item The ``error'' is the phase difference $\theta_{\mathrm{ref}} - \theta(t)$.
\item The ``proportional gain'' is the coupling strength $\kappa$.
\item The ``control signal'' is the coupling torque $\kappa \sin(\theta_{\mathrm{ref}} - \theta)$.
\item The ``integral term'' is the accumulated phase (count $M$).
\item The ``derivative term'' is the frequency difference $\dot{\theta}_{\mathrm{ref}} - \dot{\theta}$.
\end{itemize}
All these quantities are physical properties of the coupled oscillator system, as expressed in~\eqref{eq:kuramoto-coupling}. They are not computed by a separate algorithm and then applied---they \emph{are} the physics of the coupling. Therefore the coupling topology determines the control law with zero computational overhead.
\end{proof}

\begin{corollary}[Automatic Incorporation of Hard-to-Model Physics]
\label{cor:hard-physics}
Physical effects that are notoriously difficult to model in conventional architectures---tire-road friction coefficients, aerodynamic drag variations, road surface irregularities, crosswind gusts, suspension dynamics, thermal effects on tire grip---are automatically incorporated in the coupled counting loop framework. These effects need not be modeled; they are \emph{counted}.
\end{corollary}

\begin{proof}
Each physical effect manifests as a modification to some oscillator's frequency or coupling strength:
\begin{itemize}[leftmargin=*]
\item \emph{Tire-road friction}: modifies the coupling strength $\kappa_{\mathrm{tire}}$ between wheel oscillator and road surface. Low friction reduces $\kappa_{\mathrm{tire}}$, causing the wheel oscillator to decouple from the road---this is precisely what ``loss of traction'' means physically.
\item \emph{Aerodynamic drag}: modifies the damping coefficient of the vehicle body oscillator. At high speed, increased drag reduces the quality factor of the body's oscillation, which the counting loops register as a frequency shift.
\item \emph{Road irregularities}: excite the suspension oscillators at characteristic frequencies. A pothole is a sudden phase perturbation; a washboard road is a periodic forcing.
\item \emph{Crosswind}: introduces a lateral forcing oscillation on the vehicle body, which the steering counting loop registers as a phase perturbation.
\end{itemize}
By \Cref{thm:counting-identity}, the counting loops that include these oscillators automatically observe, compute, and process these effects in every tick. The SRM's triple convergence check (\Cref{def:triple-convergence}) then evaluates whether the transition remains sufficient \emph{given these actual conditions}, not given a model of these conditions. The physics is the computation.
\end{proof}

\begin{figure}[t]
\centering
\fbox{\parbox{0.9\columnwidth}{\centering\vspace{3cm}\textbf{Figure 4: Completion Morphism as Phase Transition}\\\vspace{0.3cm}Upper panel: phase portrait of the coupled oscillator network showing the transition from $\Spen$ (left attractor basin, blue) to $\Sfinal$ (right attractor basin, green). The morphism $\varphi$ traces the trajectory through the saddle point between basins (red curve). Lower three panels: longitudinal ($\Delta\St$), lateral ($\Delta\Sk$), and energy ($\Delta\Se$) morphism types with their dominant oscillator couplings highlighted. Engine--wheel coupling for longitudinal, steering--wheel coupling for lateral, transmission coupling for energy.\vspace{3cm}}}
\caption{The completion morphism as a phase transition in the coupled oscillator network. The system transitions from the phase-locked state encoding $\Spen$ to the state encoding $\Sfinal$ through transient desynchronization and resynchronization. No trajectory is computed---the coupling topology determines the transition path. The three morphism types correspond to the three S-entropy coordinates.}
\label{fig:completion-morphism}
\end{figure}

\textbf{Formal Interface Contract.}
\begin{itemize}[leftmargin=*]
\item \emph{Precondition}: SRM returns \textsc{Proceed} or \textsc{Slow}; $\Spen$ and $\Sfinal$ from PNE. If \textsc{Slow}, the morphism is executed with reduced coupling strengths (smaller $\kappa$, slower transition).
\item \emph{Postcondition}: Vehicle state transitions from $\Spen$ to $\Sfinal$ within one morphism period $\tau_\varphi \leq 2\pi/\omega_{\min}$, where $\omega_{\min}$ is the lowest relevant oscillator frequency.
\item \emph{Invariant}: During transition, the categorical distance to destination is monotonically non-increasing: $\dcat(\Sentropy(t), \Sfinal) \leq \dcat(\Spen, \Sfinal)$ for all $t$ during the morphism. Violation triggers TEM alert.
\end{itemize}


% ===================================================================
\section{Triple Equivalence Monitor}
\label{sec:tem}

The Triple Equivalence Monitor (TEM) is the watchdog subsystem that continuously verifies the fundamental identity~\eqref{eq:fundamental-identity} holds across all counting loops in the vehicle's oscillator network. It is the last line of defense: even if all other subsystems malfunction, the TEM detects inconsistency and triggers safe shutdown.

\begin{definition}[Triple Equivalence Violation]
\label{def:tev}
A \emph{triple equivalence violation} (TEV) occurs when, for any counting loop $\mathcal{L}_v$ in the network, the rate equation~\eqref{eq:rate-equation} fails:
\begin{equation}
\frac{dM_v}{dt} \neq \frac{\omega_v}{2\pi / M_v} \quad \text{or} \quad \frac{\omega_v}{2\pi / M_v} \neq \frac{1}{\langle \taup^{(v)} \rangle}.
\label{eq:tev-condition}
\end{equation}
The violation is quantified by the \emph{TEV magnitude}:
\begin{equation}
\Delta_v = \max\left(\left|\frac{dM_v}{dt} - \frac{\omega_v M_v}{2\pi}\right|, \left|\frac{\omega_v M_v}{2\pi} - \frac{1}{\langle \taup^{(v)} \rangle}\right|\right).
\end{equation}
\end{definition}

\begin{theorem}[Triple Equivalence Monitoring]
\label{thm:tem-monitoring}
If the rate equation~\eqref{eq:rate-equation} is violated for any counting loop $\mathcal{L}_v$ (i.e., $\Delta_v > \Delta_{\mathrm{threshold}}$), the system is in an inconsistent state and must trigger a safe shutdown.
\end{theorem}

\begin{proof}
The rate equation $dM/dt = \omega/(2\pi/M) = 1/\langle \taup \rangle$ is the quantitative expression of the triple identity $O \equiv C \equiv P$ (\Cref{thm:rate-equation}). Each equality encodes a different aspect:
\begin{itemize}[leftmargin=*]
\item $dM/dt = \omega M/(2\pi)$: the counting rate (computation) equals the oscillator frequency scaled by count (observation).
\item $\omega M/(2\pi) = 1/\langle \taup \rangle$: the observation-based rate equals the partition crossing rate (physical processing).
\end{itemize}

If $dM/dt \neq \omega M/(2\pi)$, then the count is incrementing at a rate inconsistent with the oscillator's frequency. This means the counting loop is either (a) missing ticks (hardware failure), (b) double-counting (interference), or (c) the oscillator frequency has changed without corresponding count rate adjustment (coupling failure). In any case, the computational state ($M$) disagrees with the observed state ($\omega$), and any decision based on this counting loop would use inconsistent information.

If $\omega M/(2\pi) \neq 1/\langle \taup \rangle$, then the oscillator-based measurement disagrees with the partition-based measurement. The physical processing (partition crossings) is occurring at a rate inconsistent with the observation (oscillator frequency). This indicates a decoupling between the computational model and physical reality.

In either case, the fundamental identity is broken for this counting loop. Since the entire architecture depends on the identity (\Cref{thm:counting-identity}), the system cannot make reliable decisions. By \Cref{thm:non-convergence}, the only safe response is to halt.
\end{proof}

\begin{proposition}[Zero-Overhead Monitoring]
\label{prop:zero-overhead}
The TEM operates with zero additional computational overhead.
\end{proposition}

\begin{proof}[Justification]
Each counting loop's tick is self-verifying. The tick event itself carries three pieces of information: (1) the time it occurred (from the oscillator), (2) the count it incremented (the computation), and (3) the partition cell it crossed (the physical state). If all three are consistent, the rate equation is satisfied. If any are inconsistent, the tick is anomalous. The ``verification'' is not a separate computation---it is the tick itself. A tick that satisfies the rate equation is a valid tick; a tick that violates it is an invalid tick. The distinction is made by the physics of the counting loop, not by a monitoring algorithm.

In practice, the TEM maintains running estimates of $dM/dt$, $\omega M/(2\pi)$, and $1/\langle \taup \rangle$ for each counting loop, which are $\BigO(|\mathcal{V}|)$ running averages updated incrementally with $\BigO(1)$ work per tick per loop. This is negligible compared to the tick rate itself.
\end{proof}

\textbf{Formal Interface Contract.}
\begin{itemize}[leftmargin=*]
\item \emph{Precondition}: All counting loops in $\mathcal{G}$ are operational (ticking).
\item \emph{Postcondition}: Any violation $\Delta_v > \Delta_{\mathrm{threshold}}$ is detected within one tick period $2\pi/\omega_v$ of the violating counting loop.
\item \emph{Invariant}: No \textsc{Proceed} decision is issued by the SRM while any TEV is active. This invariant is enforced by hardware interlock: the TEM can veto SRM output independently of software.
\end{itemize}


% ===================================================================
\section{Safety Guarantees}
\label{sec:safety}

We now collect and prove the comprehensive safety guarantees of the trajectory completion architecture. These guarantees form a defense-in-depth: each is independent and any one suffices to prevent harm.

\subsection{Partition Boundedness}

\begin{theorem}[Partition Boundedness]
\label{thm:partition-bounded}
For any state $\Sentropy = (\Sk, \St, \Se) \in [0,1]^3$ and any other state $\Sentropy' \in [0,1]^3$, the categorical distance satisfies:
\begin{equation}
\dcat(\Sentropy, \Sentropy') \leq \sqrt{3}.
\end{equation}
The system cannot diverge: the maximum categorical distance is bounded by $\sqrt{3}$ for all time, for all states, unconditionally.
\end{theorem}

\begin{proof}
The S-entropy space is the unit cube $[0,1]^3$ (\Cref{def:s-entropy}). The categorical distance $\dcat$ is bounded above by the Euclidean distance in S-entropy space~\cite{sachikonye2026trajectory}:
\[
\dcat(\Sentropy, \Sentropy') \leq \|\Sentropy - \Sentropy'\|_2.
\]
The maximum Euclidean distance in the unit cube is the space diagonal:
\[
\max_{\Sentropy, \Sentropy' \in [0,1]^3} \|\Sentropy - \Sentropy'\|_2 = \sqrt{(1-0)^2 + (1-0)^2 + (1-0)^2} = \sqrt{3}.
\]
This bound is independent of the dynamics, the counting loops, the couplings, or any other system parameter. It is a consequence of the S-entropy normalization to $[0,1]$.

The contrast with forward simulation is stark. In the conventional architecture, prediction errors satisfy $\varepsilon(t) \geq \varepsilon_0 e^{\lambda t}$ (\Cref{thm:lyapunov}), which is unbounded as $t \to \infty$. The trajectory completion architecture operates in a state space where divergence is geometrically impossible.
\end{proof}

\subsection{G\"{o}delian Safety}

\begin{theorem}[G\"{o}delian Safety]
\label{thm:godelian-safety}
The G\"{o}delian residue $\varepsilon > 0$ (the irreducible uncertainty in self-referential measurement) ensures that the system can always distinguish ``sufficient'' from ``insufficient.'' Specifically, for any counting loop $\mathcal{L}$, the measurement resolution satisfies $\delta_{\mathcal{L}} > 0$, and the sufficiency criterion is well-defined.
\end{theorem}

\begin{proof}
By G\"{o}del's incompleteness theorems~\cite{godel1931}, no sufficiently powerful formal system can fully characterize its own consistency from within. Applied to physical systems via the counting loop framework~\cite{sachikonye2026trajectory}: a counting loop measuring its own partition cell has an irreducible uncertainty $\varepsilon > 0$ because it cannot count its own counting (the observer is part of the observed system).

This is a \emph{safety feature}, not a limitation:
\begin{enumerate}
\item $\varepsilon > 0$ means there is always a minimum detectable gap $\varepsilon_{\min} > 0$. The SRM can always distinguish $\bar{\varepsilon} > \varepsilon_{\min}$ (sufficient) from $\bar{\varepsilon} < -\varepsilon_{\min}$ (insufficient).
\item If $\varepsilon = 0$ were achievable, the system could reach a state of perfect self-knowledge where $\bar{\varepsilon} = 0$ exactly, and the sufficiency criterion would be undecidable (the boundary between sufficient and insufficient would be a set of measure zero).
\item With $\varepsilon > 0$, the boundary is ``thickened'' to a band of width $\sim \varepsilon$. States within this band are mapped to \textsc{Slow} by the SRM, not to \textsc{Proceed}---the uncertainty is handled conservatively.
\end{enumerate}
The triple convergence mechanism (\Cref{def:triple-convergence}) exploits the G\"{o}delian residue: three independent measurements, each with its own $\varepsilon_i > 0$, converge only when measuring a real physical quantity. The residues $\varepsilon_i$ prevent false convergence.
\end{proof}

\subsection{Thermodynamic Security}

\begin{theorem}[Thermodynamic Security]
\label{thm:thermo-security}
Unauthorized state transitions produce detectable heating. Specifically, if a transition $\Sentropy \to \Sentropy'$ is unauthorized (not sanctioned by the SRM via triple convergence), then
\begin{equation}
\frac{d\Se}{dt} > 0 \quad \text{(heating)}.
\end{equation}
Detection requires zero computational overhead: temperature monitoring suffices.
\end{theorem}

\begin{proof}
The proof rests on the fundamental distinction between categorical and physical operations~\cite{landauer1961,bennett1973,bennett1982,sachikonye2026trajectory}.

\textbf{Authorized operations} are \emph{categorical morphisms}: transitions that correspond to valid arrows in the partition category, sanctioned by the SRM. These satisfy the commutation relation
\begin{equation}
[\Ocat, \Ophys] = 0
\label{eq:commutation}
\end{equation}
where $\Ocat$ is the categorical (logical) operator and $\Ophys$ is the physical (thermodynamic) operator. When these commute, the logical operation can be performed without disturbing the physical state, and vice versa. By Bennett's theorem~\cite{bennett1973}, logically reversible operations can in principle be performed with zero energy dissipation. Authorized categorical morphisms are logically reversible (every morphism has an inverse in the partition category~\cite{awodey2010,maclane1998}), so they satisfy $d\Se/dt = 0$ in the ideal case, and $d\Se/dt \approx 0$ in practice.

\textbf{Unauthorized operations} violate~\eqref{eq:commutation}: they are physical operations that are not valid categorical morphisms. An unauthorized transition must override the SRM's decision, which requires \emph{erasing} the SRM's state (replacing \textsc{Stop} with \textsc{Proceed} in the decision register). By Landauer's principle~\cite{landauer1961}, erasing one bit of information dissipates at least
\begin{equation}
\Delta Q \geq k_B T \ln 2 \approx 2.9 \times 10^{-21} \text{ J at 300 K}
\end{equation}
per bit. An unauthorized transition that overrides $n_{\mathrm{bits}}$ of SRM state produces heating
\begin{equation}
\Delta Q_{\mathrm{unauth}} \geq n_{\mathrm{bits}} \cdot k_B T \ln 2 > 0.
\end{equation}
This heat increases the energetic entropy $\Se$ of the affected component (processor, actuator, or communication bus). Even for $n_{\mathrm{bits}} = 1$, the heating is nonzero in principle. In practice, any realistic cyber attack involves modifying $n_{\mathrm{bits}} \gg 1$ bits (overwriting control registers, injecting false sensor data, etc.), producing measurable heating.

Temperature sensors already present on every automotive MCU, every power transistor, and every battery cell can detect anomalous $d\Se/dt > 0$. No additional computation is required---the laws of thermodynamics serve as the intrusion detection system.
\end{proof}

\begin{corollary}[Intrusion Detection via Physics]
\label{cor:intrusion-detection}
A cyber-physical attack on the vehicle---spoofed sensor data, hijacked actuator commands, falsified GPS signals, compromised CAN bus messages---necessarily produces anomalous heating at the point of injection. The TEM detects this as a TEV ($d\Se/dt > 0$ without a corresponding authorized morphism) and triggers safe shutdown. The security guarantee is physical, not cryptographic, and cannot be bypassed by software~\cite{szilard1929,leff1990}.
\end{corollary}

\subsection{Comprehensive Safety Summary}

\begin{theorem}[Comprehensive Safety]
\label{thm:comprehensive-safety}
The trajectory completion architecture satisfies the following five independent safety properties simultaneously:
\begin{enumerate}
\item \textbf{Bounded state space} (\Cref{thm:partition-bounded}): $\dcat \leq \sqrt{3}$ for all time, preventing divergence.
\item \textbf{Recognizable sufficiency} (\Cref{thm:srm-false-positive}): False-positive rate $P_{\mathrm{fp}} \leq (\delta/\varepsilon_{\min})^2 \sim 10^{-12}$.
\item \textbf{Safe default} (\Cref{thm:non-convergence}): Non-convergence $\Rightarrow$ \textsc{Stop}, proved thermodynamically safe.
\item \textbf{Consistency monitoring} (\Cref{thm:tem-monitoring}): TEV detected within one tick period.
\item \textbf{Thermodynamic security} (\Cref{thm:thermo-security}): Unauthorized transitions detected via $d\Se/dt > 0$.
\end{enumerate}
The properties are independent: each holds regardless of whether the others are operational. They form a defense-in-depth with five independent layers.
\end{theorem}

\begin{proof}
Properties 1--5 are proved in \Cref{thm:partition-bounded}, \Cref{thm:srm-false-positive}, \Cref{thm:non-convergence}, \Cref{thm:tem-monitoring}, and \Cref{thm:thermo-security} respectively. Independence: Property~1 is geometric (depends only on $[0,1]^3$ normalization). Property~2 is statistical (depends on Observer independence). Property~3 is thermodynamic (depends on second law). Property~4 is algebraic (depends on rate equation). Property~5 is information-theoretic (depends on Landauer's principle). No property's proof references another property. Therefore failure of any one property does not compromise the remaining four.

The probability that all five layers fail simultaneously is bounded by the product of individual failure probabilities. Even with conservative individual bounds of $10^{-3}$ per layer, the joint failure probability is $10^{-15}$, corresponding to one failure per $10^{15}$ operating hours---approximately $10^{11}$ years, far exceeding the lifetime of any vehicle or fleet.
\end{proof}


% ===================================================================
\section{Inter-Vehicle Coordination}
\label{sec:coordination}

We now extend the framework from single-vehicle to multi-vehicle, showing that traffic coordination emerges naturally from the categorical architecture without explicit vehicle-to-vehicle (V2V) communication protocols.

\subsection{Traffic as Categorical Gas}

\begin{theorem}[Traffic Equation of State]
\label{thm:traffic-eos}
A system of $N$ vehicles on a road network of effective volume $V_{\mathrm{road}}$ with mean flow temperature $T_{\mathrm{flow}}$ satisfies an equation of state analogous to the ideal gas law~\cite{boltzmann1896}:
\begin{equation}
P_{\mathrm{traffic}} \cdot V_{\mathrm{road}} = N \cdot k_B \cdot T_{\mathrm{flow}}
\label{eq:traffic-eos}
\end{equation}
where $P_{\mathrm{traffic}}$ is the traffic pressure (momentum flux per unit road cross-section) and $T_{\mathrm{flow}} = m \langle v^2 \rangle / k_B$ is the kinetic temperature proportional to mean-square velocity~\cite{helbing2001,lighthill1955}.
\end{theorem}

\begin{proof}
The derivation parallels the kinetic theory of ideal gases~\cite{boltzmann1896,sachikonye2026gas}. Each vehicle is treated as a ``molecule'' with mass $m$ and velocity $v$ along a one-dimensional road. The mean kinetic energy per vehicle defines the temperature: $\frac{1}{2}m\langle v^2 \rangle = \frac{1}{2}k_B T_{\mathrm{flow}}$.

The traffic pressure is the momentum transfer rate per unit cross-section at a road boundary (on-ramp, traffic light, toll booth). Consider a road segment of length $L$ and cross-section $A$ containing $N$ vehicles. The pressure exerted on the downstream boundary is
\[
P_{\mathrm{traffic}} = \frac{N m \langle v^2 \rangle}{V_{\mathrm{road}}} = \frac{N m \langle v^2 \rangle}{A \cdot L}
\]
by the standard kinetic theory derivation (momentum $mv$ transferred at rate $v/L$ per vehicle, summed over $N$ vehicles). Since $m\langle v^2 \rangle = k_B T_{\mathrm{flow}}$:
\[
P_{\mathrm{traffic}} = \frac{N k_B T_{\mathrm{flow}}}{V_{\mathrm{road}}}.
\]
Rearranging gives $P_{\mathrm{traffic}} \cdot V_{\mathrm{road}} = N k_B T_{\mathrm{flow}}$~\cite{sachikonye2026equations}.
\end{proof}

\begin{definition}[Traffic Phase Transitions]
\label{def:traffic-phases}
The categorical gas of $N$ vehicles exhibits three thermodynamic phases, characterized by the energetic entropy $\Se$:
\begin{itemize}[leftmargin=*]
\item \textbf{Gas phase (free flow)}: $\Se \approx 0$. Vehicles are independent counting loop networks with negligible inter-vehicle coupling. Each vehicle occupies its own partition cells exclusively. Mean free path $\ell \gg$ vehicle length. Velocity distribution is approximately Maxwellian.

\item \textbf{Liquid phase (coordinated flow)}: $0 < \Se < \Se^{\mathrm{crit}}$. Partial coupling between vehicles through shared partition cells. Vehicles adjust speeds to maintain flow coherence (``traffic waves'' emerge). Mean free path $\ell \sim$ vehicle length. Velocity distribution narrows around the mean.

\item \textbf{Crystal phase (platoon)}: $\Se > \Se^{\mathrm{crit}}$. Phase-locked inter-vehicle coupling. Vehicles form rigid platoons with fixed spacing $d_0$, moving as a single oscillatory entity. Mean free path $\ell \to 0$ within the platoon. Velocity variance $\to 0$.
\end{itemize}
The critical energetic entropy $\Se^{\mathrm{crit}}$ is determined by the critical coupling strength of the Kuramoto model: $\Se^{\mathrm{crit}} = k_B T \ln(2\pi g(0)/\kappa_c)$ where $g(0)$ is the natural frequency distribution at zero detuning~\cite{kuramoto1984}.
\end{definition}

\begin{theorem}[Emergent Coordination]
\label{thm:emergent-coordination}
Phase transitions between traffic states (gas $\leftrightarrow$ liquid $\leftrightarrow$ crystal) occur without explicit communication between vehicles. The coordination emerges from the shared partition space.
\end{theorem}

\begin{proof}
Each vehicle's counting loops observe the partition cells they occupy (\Cref{def:counting-loop}). When two vehicles $A$ and $B$ are in adjacent or overlapping partition cells (i.e., they are physically close on the road), their counting loops automatically register the shared occupancy.

Specifically, vehicle $A$'s wheel speed counting loop observes the partition cell of the road ahead. If vehicle $B$ occupies that cell, $B$'s physical presence modifies the cell's state (e.g., vehicle $B$ reflects radar signals, occludes visual field, or physically constrains $A$'s trajectory). Vehicle $A$'s counting loops register these modifications as changes in their tick patterns. This is observation ($O$), which by \Cref{thm:counting-identity} is also computation ($C$) and processing ($P$). No data is ``transmitted'' from $B$ to $A$---$A$ simply counts the state of the partition cell, which includes $B$'s presence.

The inter-vehicle coupling is therefore mediated by the partition space itself. As density increases (more vehicles per partition cell), the coupling strength increases because each vehicle's counting loops are more strongly influenced by neighbors. By Kuramoto's synchronization threshold~\cite{kuramoto1984}, there exists a critical density $\rho_c$ at which spontaneous velocity synchronization occurs (gas $\to$ liquid transition). At higher density $\rho_{c2} > \rho_c$, full phase-locking produces platoons (liquid $\to$ crystal transition).

No V2V communication protocol (DSRC, C-V2X, or otherwise) is involved. The ``communication medium'' is the physical space that the vehicles share---the same partition space they use for navigation~\cite{sachikonye2026sango}.
\end{proof}

\begin{corollary}[No V2V Protocol Required]
\label{cor:no-v2v}
Dedicated Short-Range Communication (DSRC), Cellular V2X (C-V2X), or any other explicit inter-vehicle communication protocol is unnecessary in the categorical architecture. Vehicles coordinate through the physics of shared space occupation, which is counted by their existing sensor counting loops. This eliminates the security vulnerabilities, bandwidth limitations, and standardization challenges of explicit V2V protocols.
\end{corollary}

\begin{proposition}[Congestion Prediction from Equation of State]
\label{prop:congestion}
Traffic congestion corresponds to a phase transition from the gas or liquid phase to a degenerate state where $P_{\mathrm{traffic}} \to \infty$ at fixed $V_{\mathrm{road}}$ (or equivalently, effective $V_{\mathrm{road}} \to 0$ as available road capacity is consumed). From~\eqref{eq:traffic-eos}, this occurs when $N \cdot T_{\mathrm{flow}} / V_{\mathrm{road}}$ exceeds a critical threshold. Since $N$ (number of nearby vehicles) and $T_{\mathrm{flow}}$ (mean-square velocity) are counted by each vehicle's Observer counting loops, incipient congestion is detectable locally without any centralized traffic monitoring system.
\end{proposition}


% ===================================================================
\section{Comparison with Conventional Architectures}
\label{sec:comparison}

We now provide a systematic comparison between the trajectory completion architecture and conventional autonomous driving architectures across four dimensions: complexity, information requirements, safety guarantees, and failure modes.

\subsection{Complexity Comparison}

\begin{table}[t]
\centering
\caption{Computational complexity comparison. $n$: agents, $T$: horizon, $\Delta t$: time step, $N$: road segments, $K$: Observers, $V$: intersections, $E$: road edges, $d$: state dimension.}
\label{tab:complexity}
\small
\begin{tabular}{@{}lcc@{}}
\toprule
\textbf{Operation} & \textbf{Conventional} & \textbf{Traj.\ Completion} \\
\midrule
Perception & $\BigO(n_{\mathrm{pts}} \cdot d)$ & $\BigO(K)$ \\
Prediction & $\BigO(n^2 T / \Delta t)$ & --- (eliminated) \\
Path planning & $\BigO(E + V \log V)$ & $\BigO(\log_3 N)$ \\
Control & $\BigO(d^3)$ (MPC) & $\BigO(1)$ (morphism) \\
Safety check & $\BigO(n^2)$ (collision) & $\BigO(K)$ (convergence) \\
V2V comm. & $\BigO(N_{\mathrm{msg}} \cdot B)$ & --- (eliminated) \\
\midrule
\textbf{Total per cycle} & $\BigO(n^2 T / \Delta t)$ & $\BigO(\log_3 N + K)$ \\
\bottomrule
\end{tabular}
\end{table}

\Cref{tab:complexity} summarizes the complexity comparison. The trajectory completion architecture achieves an exponential reduction in computational complexity through three mechanisms:
\begin{itemize}[leftmargin=*]
\item \textbf{Prediction eliminated}: The dominant $\BigO(n^2 T/\Delta t)$ cost is removed entirely.
\item \textbf{Logarithmic navigation}: $\BigO(\log_3 N)$ vs $\BigO(E + V \log V)$---exploiting tree structure vs.\ treating the road as an unstructured graph.
\item \textbf{Implicit control}: $\BigO(1)$ morphism execution through physical coupling vs.\ $\BigO(d^3)$ model-predictive control requiring matrix inversion at each step.
\end{itemize}

\subsection{Information-Theoretic Comparison}

\begin{table}[t]
\centering
\caption{Information requirements comparison.}
\label{tab:information}
\small
\begin{tabular}{@{}lcc@{}}
\toprule
\textbf{Quantity} & \textbf{Conventional} & \textbf{Traj.\ Completion} \\
\midrule
Sensor data processed & $10^8$--$10^9$ bits/frame & $2$--$5$ bits/decision \\
Prediction model size & $10^6$--$10^9$ parameters & 0 parameters \\
Map representation & $\BigO(V + E)$ graph & $\BigO(D)$ tree depth \\
Control parameters & $10^2$--$10^4$ & 0 (implicit in coupling) \\
V2V messages & $\BigO(N)$ per cycle & 0 \\
\bottomrule
\end{tabular}
\end{table}

The information-theoretic comparison (\Cref{tab:information}) reflects \Cref{thm:info-waste}: the trajectory completion architecture processes only the mutual information $I(D; A_Q) \approx 2$--$5$ bits per decision, while conventional architectures process the full environment entropy $H(D) \approx 10^8$--$10^9$ bits. This $10^3$--$10^7$-fold reduction in information processing explains the $10^5$-fold reduction in computational complexity.

\subsection{Safety Guarantee Comparison}

\begin{table}[t]
\centering
\caption{Safety guarantee comparison.}
\label{tab:safety}
\small
\begin{tabular}{@{}lcc@{}}
\toprule
\textbf{Property} & \textbf{Conventional} & \textbf{Traj.\ Completion} \\
\midrule
State bounded & No ($\varepsilon \to \infty$) & Yes ($\dcat \leq \sqrt{3}$) \\
Failure prob. & $P_f \to 1$ (\Cref{thm:unbounded-failure}) & $P_f \leq (\delta/\varepsilon)^2$ \\
Default action & Undefined & \textsc{Stop} (proved safe) \\
Intrusion detect. & Software (bypassable) & Physics ($d\Se/dt$) \\
Consistency & Ad hoc watchdogs & Triple equivalence \\
\# safety layers & 1--2 (redundancy) & 5 (independent) \\
\bottomrule
\end{tabular}
\end{table}

\subsection{Failure Mode Comparison}

The failure modes of the two architectures are qualitatively different:

\textbf{Conventional architecture: silent failure via prediction error.} The predicted state diverges from reality (\Cref{thm:lyapunov}). The system does not know its prediction is wrong until the error manifests physically (e.g., collision). By \Cref{thm:unbounded-failure}, such failures are mathematically inevitable over sufficient operating time. The failure is \emph{silent}: the system has full confidence in its (wrong) prediction.

\textbf{Trajectory completion architecture: loud failure via non-convergence.} The three perspectives disagree (\Cref{def:triple-convergence} violated). The system knows immediately that something is wrong---the non-convergence \emph{is} the detection. The SRM commands \textsc{Stop} (\Cref{thm:non-convergence}), which is proved safe. The failure is \emph{loud}: the system explicitly flags its inability to confirm sufficiency.

The architectural transformation converts \emph{silent} failures (dangerous, undetectable) into \emph{loud} failures (safe, immediately handled). This is the fundamental safety advantage.

\begin{figure}[t]
\centering
\fbox{\parbox{0.9\columnwidth}{\centering\vspace{2.5cm}\textbf{Figure 5: Failure Mode Comparison}\\\vspace{0.3cm}Left panel: Conventional architecture---prediction error $\varepsilon(t)$ (red curve) grows exponentially (\Cref{thm:lyapunov}), crossing the safety threshold (dashed line) at $T_{\mathrm{safe}}$. Gray region above threshold is unsafe; the system has no mechanism to detect entry into this region. Right panel: Trajectory completion architecture---non-convergence magnitude $|\varepsilon_{\mathrm{osc}} - \varepsilon_{\mathrm{cat}}|$ (blue curve) is monitored continuously. When it exceeds $\delta$ (dashed line), the system immediately commands \textsc{Stop}. The unsafe region is never entered.\vspace{2.5cm}}}
\caption{Failure mode comparison. Conventional architectures fail silently when exponentially growing prediction error crosses the safety threshold---the system cannot detect this crossing. The trajectory completion architecture monitors non-convergence continuously and triggers a provably safe stop the moment agreement breaks down.}
\label{fig:failure-comparison}
\end{figure}


% ===================================================================
\section{Implementation}
\label{sec:implementation}

\subsection{Mapping to Vehicle Hardware}

The trajectory completion architecture does not require novel hardware. Existing vehicle components already implement counting loops:

\begin{itemize}[leftmargin=*]
\item \textbf{Wheel speed sensors} (ABS reluctor ring + Hall effect sensor): counting loop with $\omega_{\mathrm{wheel}}$, partition cell = wheel angular position. Existing on every ABS-equipped vehicle since $\sim$1990.
\item \textbf{Crankshaft/camshaft position sensors}: counting loop with $\omega_{\mathrm{eng}}$, partition cell = engine rotational state. Existing on every fuel-injected vehicle.
\item \textbf{Steering angle sensor}: counting loop with $\omega_{\mathrm{steer}}$, partition cell = steering column angular position. Existing on every ESC-equipped vehicle since $\sim$2012.
\item \textbf{IMU (inertial measurement unit)}: MEMS oscillator counting loop with $\omega_{\mathrm{IMU}} \sim 10^3$--$10^4$ Hz, partition cell = acceleration/angular rate vector. Standard on modern vehicles.
\item \textbf{Camera frame clock}: counting loop with $\omega_{\mathrm{cam}} \sim 30$--$60$ Hz, partition cell = image region. Standard on vehicles with ADAS.
\item \textbf{LiDAR rotation motor}: counting loop with $\omega_{\mathrm{LiDAR}} \sim 10$--$20$ Hz, partition cell = azimuthal sector. Present on vehicles with LiDAR.
\item \textbf{MCU clock crystal}: counting loop with $\omega_{\mathrm{MCU}} \sim 10^8$--$10^9$ Hz, partition cell = processor register state. Present on every vehicle ECU.
\end{itemize}

The key insight is that these counting loops already exist, already operate, and already produce the tick streams needed by the architecture. They are simply not recognized as elements of a unified computing framework. The trajectory completion architecture unifies them into a single coupled counting loop network (\Cref{thm:identity-unification}).

\subsection{Mapping to the Verum Codebase}

The architecture maps directly to the existing verum implementation~\cite{sachikonye2026buhera,sachikonye2026vahera}:

\begin{itemize}[leftmargin=*]
\item \texttt{sighthound} module $\to$ Observer counting loops: implements sensor abstraction as categorical observation, providing raw counting loop ticks to the CSM.
\item \texttt{gusheshe} module $\to$ Sufficiency Recognition Module: evaluates state transitions for sufficiency via the triple convergence check of \Cref{alg:sufficiency}.
\item \texttt{verum-core} module $\to$ Penultimate Navigation Engine: implements backward navigation through the partition tree (\Cref{alg:backward-nav}).
\item \texttt{vaHera} interface~\cite{sachikonye2026vahera,sachikonye2026zangalewa}: the user declares a destination via the vaHera categorical scripting language. The system resolves the destination to a partition tree leaf $\Sfinal$ and initiates the trajectory completion cycle: PNE computes $\Spen$, SRM evaluates sufficiency, CME executes morphisms, TEM monitors throughout.
\end{itemize}

The user interaction is minimal by design. The driver (or passenger) declares an intent---a destination---and the system completes the trajectory. There are no intermediate waypoints, no route selection, no speed settings. The system determines all of these through backward navigation and sufficiency recognition, just as a human driver decides the specific lane changes and speed adjustments without conscious route micro-management.

\subsection{Computational Requirements}

\begin{proposition}[Automotive MCU Sufficiency]
\label{prop:mcu}
The trajectory completion architecture executes on existing automotive microcontrollers (e.g., Infineon AURIX TC3xx at 300 MHz, NXP S32G at 1 GHz) without requiring GPUs, TPUs, or dedicated AI accelerators.
\end{proposition}

\begin{proof}[Justification]
The computational load per decision cycle is:
\begin{enumerate}
\item \textbf{CSM} (\Cref{def:csm-estimation}): weighted average of $K \approx 20$ Observer S-entropy estimates. Each average requires $3$ multiplications and $3$ additions (one per S-entropy coordinate). Total: $6K = 120$ floating-point operations plus $K = 20$ variance updates ($\sim 10$ ops each): $\sim 320$ ops.
\item \textbf{PNE} (\Cref{alg:backward-nav}): tree traversal of depth $D \leq 15$ with $3$ comparisons per level: $\sim 45$ comparisons $= \sim 90$ ops.
\item \textbf{SRM} (\Cref{alg:sufficiency}): three gap computations ($\sim 20$ ops each) and three comparisons: $\sim 65$ ops.
\item \textbf{TEM} (\Cref{def:tev}): rate equation verification for $|\mathcal{V}| \approx 50$ counting loops, $3$ ops each: $\sim 150$ ops.
\end{enumerate}
Total: $\sim 625$ operations per decision cycle. At a decision rate of 100 Hz (one decision per 10 ms, matching the fastest ADAS cycle time): $62{,}500$ operations/second.

An Infineon AURIX TC3xx at 300 MHz executes approximately $300 \times 10^6$ integer operations per second (with three cores). The trajectory completion architecture uses $6.25 \times 10^4 / 3 \times 10^8 \approx 0.02\%$ of the available computation. Even accounting for overhead (memory access, branching, OS scheduling), the utilization is below $1\%$.

This extreme efficiency is a direct consequence of \Cref{thm:info-waste}: the architecture processes $2$--$5$ bits per decision, not $10^8$--$10^9$ bits. The seven-to-nine orders of magnitude reduction in information processing translates directly to computational savings. GPU-class hardware, which is required by conventional architectures to process the full sensor entropy in real time, is entirely unnecessary.
\end{proof}


% ===================================================================
\section{Conclusion}
\label{sec:conclusion}

We have presented a complete computing architecture for autonomous vehicles that replaces the conventional perception--prediction--planning--control pipeline with a unified framework based on trajectory completion in bounded phase space. The key results are:

\begin{enumerate}[leftmargin=*]
\item \textbf{Forward simulation is provably insufficient} (\Cref{thm:lyapunov,thm:unbounded-failure,thm:info-waste}): Lyapunov divergence, unbounded failure probability, and information-theoretic waste establish that prediction-based architectures cannot achieve the safety levels required for autonomous driving.

\item \textbf{The counting loop is a sufficient primitive} (\Cref{thm:counting-identity,thm:rate-equation,cor:no-other-primitive}): A single computational element---the counting loop---simultaneously performs observation, computation, and processing through the fundamental identity $O(x) \equiv C(x) \equiv P(x)$.

\item \textbf{Navigation is logarithmic} (\Cref{thm:nav-complexity,thm:nav-correctness}): Backward navigation through the ternary partition tree finds the unique penultimate state in $\BigO(\log_3 N)$ time, replacing $\BigO(E + V \log V)$ conventional path planning.

\item \textbf{Safety is provable and layered} (\Cref{thm:comprehensive-safety}): Five independent safety properties---partition boundedness, recognizable sufficiency, safe default, consistency monitoring, and thermodynamic security---provide defense-in-depth with formal guarantees.

\item \textbf{Coordination is emergent} (\Cref{thm:emergent-coordination,cor:no-v2v}): Inter-vehicle coordination arises as a thermodynamic phase transition in the categorical gas, requiring no explicit V2V communication protocol.

\item \textbf{Existing hardware suffices} (\Cref{prop:mcu}): The architecture requires $\sim 625$ operations per decision cycle, executable on standard automotive MCUs with $>99\%$ capacity to spare.
\end{enumerate}

The trajectory completion architecture does not incrementally improve the conventional pipeline---it replaces it with a framework whose mathematical structure matches the physical structure of driving. The conventional pipeline asks: ``What will happen?'' and fails because the answer is unknowable (\Cref{thm:lyapunov}). The trajectory completion architecture asks: ``Is there enough?'' and succeeds because the answer is always knowable through triple convergence (\Cref{thm:sufficiency}).

The deepest insight is that driving does not require prediction because humans have never used prediction to drive. What appears to be prediction (``I think that car will stay in its lane'') is actually sufficiency recognition (``there is enough space for my next action regardless of what that car does''). The trajectory completion architecture formalizes this distinction and proves its consequences.

The road to safe autonomous driving does not pass through better prediction models, larger training datasets, or more powerful GPUs. It passes through the recognition that prediction was never the right question. The right question is sufficiency, and sufficiency is recognizable in $\BigO(\log_3 N)$ time with $2$--$5$ bits of information on hardware that already exists in every car on the road.

\begin{figure}[t]
\centering
\fbox{\parbox{0.9\columnwidth}{\centering\vspace{2.5cm}\textbf{Figure 6: Paradigm Comparison Summary}\\\vspace{0.3cm}Two-panel comparison. Left: conventional pipeline with four serial stages, each a potential failure point, processing $10^8$--$10^9$ bits with $\BigO(n^2T/\Delta t)$ complexity and unbounded error growth. Right: trajectory completion architecture as a unified counting loop network, processing $2$--$5$ bits with $\BigO(\log_3 N)$ complexity, bounded error ($\dcat \leq \sqrt{3}$), five independent safety layers, and emergent multi-vehicle coordination. The conventional architecture fights against mathematical impossibility; the trajectory completion architecture works with mathematical structure.\vspace{2.5cm}}}
\caption{Paradigm summary. The conventional pipeline (left) is a serial chain of independent, failure-prone stages processing orders of magnitude more information than driving requires. The trajectory completion architecture (right) is a unified counting loop network that processes only the relevant bits, navigates in logarithmic time, provides provable safety through five independent layers, and coordinates with other vehicles through shared physics---all on existing hardware.}
\label{fig:paradigm-comparison}
\end{figure}

\section*{Acknowledgments}
This work builds on the trajectory completion computing framework and the Buhera operating system developed at the Technical University of Munich. The author thanks the verum project contributors for implementation insights.

\bibliographystyle{unsrt}
\bibliography{references}

\end{document}
